\documentclass[12pt,a4paper]{report}
\usepackage[utf8]{inputenc}
\usepackage[ngerman]{babel}
\usepackage{geometry}
\geometry{a4paper, margin=2.5cm}
\usepackage{setspace}
\onehalfspacing
\usepackage{graphicx}
\usepackage{booktabs}
\usepackage{array}
% Linksbuendige Spalte mit fester Breite (fuer Tabellen mit viel Text)
\newcolumntype{L}[1]{>{\raggedright\arraybackslash}p{#1}}
\usepackage{amsmath}
\usepackage{listings}
\usepackage{xcolor}
\usepackage{hyperref}
\usepackage{caption}
\usepackage{csquotes}
\usepackage[backend=biber, style=numeric, sorting=none]{biblatex}
\usepackage{titlesec}
\usepackage{tikz}
\usetikzlibrary{positioning,fit,calc,arrows.meta,shapes.geometric,backgrounds}
\usepackage{float}
\hypersetup{hidelinks}
\usepackage{subcaption}
\usepackage{wrapfig}
\usetikzlibrary{decorations.pathreplacing}
% ---- Kapitelüberschriften kompakt ----
\titleformat{\chapter}[hang]
  {\normalfont\LARGE\bfseries}{\thechapter\quad}{0pt}{}
\titlespacing*{\chapter}{0pt}{0pt}{20pt}
% ---- Literatur und Quellen
\addbibresource{Literatur.bib}

% ---- Gliederungstiefe: dritte Ebene nummerieren, Inhaltsverzeichnis bis Ebene zwei ----
\setcounter{secnumdepth}{2}
\setcounter{tocdepth}{2}

% ---- Code-Listing-Style (fuer Python/JS-Ausschnitte) ----
\lstdefinestyle{code}{
    basicstyle=\ttfamily\small,
    breaklines=true,
    frame=single,
    numbers=left,
    numberstyle=\tiny,
    keywordstyle=\color{blue},
    commentstyle=\color{gray},
    stringstyle=\color{teal},
    backgroundcolor=\color{gray!5}
}
\lstset{style=code}
\lstdefinelanguage{JavaScript}{
  keywords={const,let,var,function,return,if,else,for,of,in,new,this,typeof,try,catch,=>},
  keywordstyle=\color{blue}\bfseries,
  ndkeywords={true,false,null,undefined},
  comment=[l]{//}, morecomment=[s]{/*}{*/},
  morestring=[b]', morestring=[b]",
  sensitive=true
}
\renewcommand{\lstlistingname}{Quelltext}
\emergencystretch=2em

\begin{document}

% ==================== TITELSEITE ====================
\begin{titlepage}
    \centering
    \vspace*{2cm}

    {\LARGE\bfseries Ostbayerische Technische Hochschule Regensburg}\\[0.5cm]
    {\Large Fakultät Informatik und Mathematik}\\[3cm]

    {\huge\bfseries Multi-Layer Visualisierung von Tiersichtungsdaten in einer Online-Kartenplattform}\\[3cm]

    {\Large Bachelorarbeit}\\[1cm]

    vorgelegt von\\[0.5cm]
    {\Large\bfseries Yassir Sbai}\\[0.3cm]
    Matrikelnummer: 3341177\\[3cm]

    Betreuer und Prüfer: Prof. Dr. Jan Dünnweber\\


    \vfill

    Regensburg, \today
\end{titlepage}

% ==================== INHALTSVERZEICHNIS ====================
\tableofcontents
\newpage

% ====================================================================
% ====================================================================
\chapter{Einleitung}
\label{chap:einleitung}

\section{Motivation und Problemstellung}
\label{sec:motivation}

Wer frei lebende Wildtiere erforscht oder schützt, braucht heute digitale Werkzeuge, die
Bewegungsmuster in Raum und Zeit sichtbar machen. Bei großen Arten wie Bären, bei denen es
zu Konflikten mit Menschen kommen kann, fallen dabei sehr unterschiedliche Daten an:
GPS-Daten besenderter Tiere, automatisch erfasste Kamerabilder und Pfade aus
Bewegungsmodellen. Damit ein Forschungsteam damit arbeiten kann, müssen diese Daten nicht
nur gespeichert, sondern auch verständlich dargestellt werden -- im Regelfall auf einer
interaktiven Karte, auf der räumliche Zusammenhänge unmittelbar zu sehen sind.

Die beiden Datensätze dieser Arbeit sind nicht nur unterschiedlich formatiert, sie sind von
grundlegend verschiedener Art. Aus einer Langzeitstudie in den rumänischen Karpaten liegen
\emph{Trajektorien} vor: rund 60.000 GPS-Ortungen von 13 besenderten Europäischen
Braunbären (\emph{Ursus arctos}) über vier Jahre. Jede Ortung gehört zu einem bekannten
Tier und steht in zeitlicher Reihenfolge zu dessen übrigen Ortungen.

Aus der japanischen Region Fukushima liegen dagegen \emph{Detektionen} vor: zeitgestempelte
Fotoauslösungen fest installierter Wildtierkameras. Welches einzelne Tier ausgelöst hat, ist
nicht bekannt, und damit auch nicht, ob zwei Aufnahmen zum selben Tier gehören. Dieser
Unterschied ist für die Arbeit zentral, weil er bestimmt, welche Auswertungen möglich sind:
Pfade lassen sich nur auf der einen Seite vergleichen, Tagesrhythmen dagegen auf beiden.

An der Ostbayerischen Technischen Hochschule (OTH) Regensburg entstanden in einer
Kooperation mit der Universität Aizu in Japan drei eigenständige Teilsysteme. Das erste ist
eine Kartenanwendung mit zwei getrennten Seiten: eine für die Gerätestandorte und die dort
erfassten Kameraaufnahmen in der Region Aizu (Präfektur Fukushima, Honshū), eine für die
Bewegungspfade der Karpatenbären \cite{LorenzBA}. Das zweite ist ein Modell, das
Bärenbewegungen auf Grundlage der realen GPS-Daten aus den Karpaten simuliert
\cite{BreuBA}. Das dritte holt Wildtierbilder automatisch aus einem gemeinsam genutzten
E-Mail-Postfach \cite{LelchatBA}. Alle drei entstanden in getrennten studentischen Arbeiten,
und nur das erste hatte eine Benutzeroberfläche.

Daraus ergeben sich drei Einschränkungen.

Erstens ließen sich die Daten nur getrennt betrachten. Die Streamlit-Anwendung schaltete
über die Seitenleiste zwischen den beiden Seiten um, und jede erzeugte ihre eigene,
unabhängige Karte auf Basis der JavaScript-Bibliothek Leaflet, eingebunden über deren
Python-Wrapper Folium. Gerätestandorte und Bewegungspfade lagen nie auf derselben Karte.

Zweitens war die Technik nicht für die angestrebte Datenmenge gedacht. Sollen alle Bestände
gleichzeitig auf einer Karte liegen, steigt die Zahl der dargestellten Objekte stark an.
Leaflet legt für jedes Objekt ein eigenes Element im Dokumentbaum des Browsers an, und der
Aufwand wächst mit deren Zahl.

Drittens brachte jede der drei Arbeiten ihre eigene Form der Datenhaltung mit. Für sich
genommen war das jeweils angemessen, erschwert aber einen einheitlichen Zugriff.

Das Bewegungsmodell lag zudem nur als eigenständiger Quelltext vor, denn Gegenstand jener
Arbeit war das Modell selbst und nicht seine Einbindung in eine Anwendung. Die erzeugten
Pfade waren damit zwar vorhanden, aber auf keiner Karte sichtbar und nicht gegen die realen
Daten prüfbar. Genau hier setzt diese Arbeit an: Ein Forschungsteam, das reale
Bewegungsmuster mit modellierten Pfaden und mit den Standorten der Geräte in Beziehung
setzen will, kann das im Ausgangszustand nicht.

\section{Zielsetzung der Arbeit}
\label{sec:zielsetzung}

Ziel dieser Arbeit ist eine einheitliche Kartenplattform, die die drei Teilsysteme in einer
gemeinsamen, mehrschichtigen (Multi-Layer) Darstellung zusammenführt. Dazu kommen Werkzeuge,
mit denen sich simulierte Bewegungen gegen reale Beobachtungsdaten prüfen lassen. Im
Einzelnen sind das die folgenden Teilziele.

\subsection*{Integration und technische Basis}

\begin{itemize}
    \item \textbf{Technologiewechsel der Kartenvisualisierung:} Die bisherige Darstellung
    über Folium und Leaflet wird durch MapLibre GL JS abgelöst, eine quelloffene
    Kartenbibliothek, die auf WebGL und Vektorkacheln aufsetzt. Damit sollen große
    Datenmengen flüssig darstellbar werden und zusätzliche Darstellungsformen möglich sein,
    etwa Satellitenkarten und dreidimensionales Gelände.

    \item \textbf{Systemintegration:} Die drei Teilsysteme werden in einer Anwendung mit
    einer einzigen, gemeinsam genutzten Kartenkomponente zusammengeführt, auf der sich
    einzelne Datenebenen unabhängig voneinander ein- und ausblenden lassen. Bewegungsmodell
    und Bildextraktion waren über die Oberfläche bisher gar nicht erreichbar; ihre
    Einbindung ist deshalb ein wesentlicher neuer Teil dieser Arbeit.

    \item \textbf{Vereinheitlichung der Datenhaltung:} Die bisher teils lokal gehaltenen
    Daten werden in eine zentrale MariaDB-Datenbank überführt, damit mehrere Personen
    gleichzeitig damit arbeiten können und sich serverseitig filtern lässt.

    \item \textbf{Interaktions- und Filterkonzept:} Forschende sollen einzelne Bären,
    Zeiträume, Jahresphasen und Darstellungsarten auswählen können. Außerdem wird die aus
    dem Ausgangssystem bekannte Funktion wiederhergestellt, eigene Marker und Flächen auf
    der Karte anzulegen und wieder zu entfernen.
\end{itemize}

\subsection*{Analyse- und Validierungswerkzeuge}

\begin{itemize}
    \item \textbf{Bewegungsanimation:} Eine Animationskomponente macht den zeitlichen
    Verlauf einer Bewegung abspielbar, statt ihn nur als fertige Linie zu zeigen. Wichtig
    ist dabei die \emph{synchrone} Gegenüberstellung: Zu einer ausgewählten realen
    Trajektorie werden die Alternativen des Modells für denselben Bären, denselben Tag und
    dieselben Zeitpunkte angezeigt. Abweichungen zwischen Modell und Beobachtung werden so
    direkt auf der Karte sichtbar.

    \item \textbf{Übertragung des Bewegungsmodells auf die japanischen Kameradaten:} Die
    Schrittauswahl des Modells wird erweitert, damit es auch in einer dichter besiedelten
    und stärker erschlossenen Landschaft plausible Pfade erzeugt. Dazu gehören eine aus
    Daten geschätzte Richtungsträgheit, eine Gewichtung des Geländes anhand von
    OpenStreetMap-Daten und eine überarbeitete Richtungswahl.

    \item \textbf{Statistische Validierung über den Tagesrhythmus:} Da sich aus einzelnen
    Kameradetektionen keine Pfade rekonstruieren lassen, wird der Vergleich auf die Zeit
    verlagert. Über den Überlappungskoeffizienten $\Delta$ nach Ridout und Linkie
    \cite{RidoutLinkie2009} wird gemessen, wie stark die Tagesaktivität simulierter Bären
    mit der der realen japanischen Fotoauslösungen übereinstimmt, zusammen mit einer
    Unsicherheitsangabe über ein Bootstrap-Verfahren.
\end{itemize}

\subsection*{Abgrenzung des Leistungsumfangs}

Die fachlichen Modelle werden weder vollständig neu entworfen noch unverändert übernommen.
Übernommen werden die in der Vorarbeit \cite{BreuBA} aus den Karpaten-Daten abgeleiteten
Verteilungen für Geschwindigkeit und Schrittlänge. Erweitert wird die
Richtungsentscheidung: Die Richtungsträgheit wird aus Paaren aufeinanderfolgender Schritte
der realen Trajektorien geschätzt, eine Gewichtung des Geländes kommt hinzu, und die
Kandidatenrichtungen werden geschichtet über den Vollkreis gezogen, sodass ihre
Gleichverteilung schon aus der Konstruktion folgt.

Nicht Gegenstand dieser Arbeit sind die Bildverarbeitung und die automatische Artbestimmung
auf Kamerafallenbildern; ob ein Bild einen Bären zeigt, wird als gegeben vorausgesetzt.
Ebenso ist die Anwendung nicht für Echtzeitdaten ausgelegt -- eingehende Bilder werden in
Stapeln verarbeitet. Die Simulation ist schließlich kein Vorhersagewerkzeug für ein
einzelnes Tier: Sie erzeugt zufällige Pfade und wird deshalb nicht daran gemessen, ob ein
einzelner Pfad passt, sondern an statistischen Verteilungsmaßen.

\section{Aufbau der Arbeit}
\label{sec:aufbau}

Die Arbeit ist wie folgt aufgebaut.

Kapitel~\ref{chap:grundlagen} legt die Grundlagen dar. Abschnitt~\ref{sec:geodaten}
behandelt Geodaten und ihre Darstellung im Web, Abschnitt~\ref{sec:bibliotheken} die
infrage kommenden Kartenbibliotheken und ihre Einbindung in den verwendeten
Anwendungsrahmen. Abschnitt~\ref{sec:fachliche-grundlagen} führt die fachlichen Grundlagen
ein, die für die Modellerweiterung und die statistische Auswertung gebraucht werden.

Kapitel~\ref{chap:analyse} analysiert die bestehenden Systeme, ihre Datenbestände und die
übernommenen Modelle. Besonderes Gewicht liegt auf den festgestellten technischen
Einschränkungen, die mehrere Entwurfsentscheidungen in
Kapitel~\ref{chap:implementierung} begründen. Am Ende des Kapitels stehen daraus
abgeleitete, nummerierte Anforderungen.

Kapitel~\ref{chap:implementierung} beschreibt die Implementierung: die Kartenvisualisierung,
die Integration der Teilsysteme, die Vereinheitlichung der Datenhaltung, die
Bewegungsanimation, die Übertragung des Bewegungsmodells auf die japanischen Kameradaten
sowie den statistischen Vergleich der Tagesaktivität.

Kapitel~\ref{chap:diskussion} gleicht das Ergebnis gegen die Anforderungen ab, ordnet das
erzielte $\Delta$-Ergebnis ein und benennt die Grenzen der Arbeit.
Kapitel~\ref{chap:fazit} fasst zusammen und gibt einen Ausblick.

% ====================================================================
\chapter{Grundlagen}
\label{chap:grundlagen}

Dieses Kapitel stellt die Grundlagen bereit, auf denen die spätere Umsetzung aufbaut. Die
Abschnitte~\ref{sec:geodaten} und~\ref{sec:bibliotheken} behandeln die technische Seite:
wie Geodaten beschrieben werden, wie Web-Karten sie darstellen und in welchem Rahmen die
Anwendung läuft. Abschnitt~\ref{sec:fachliche-grundlagen} behandelt die fachliche Seite:
wie Tierbewegung modelliert wird und wie sich aus Kamerafallendaten Aussagen gewinnen
lassen.

% =====================================================================
\section{Geodaten und Web-Kartenvisualisierung}
\label{sec:geodaten}

\subsection{Vektor- und Rasterdaten}

Geodaten sind Daten mit Raumbezug: Jeder Datensatz gehört zu einer Position auf der
Erdoberfläche. Es gibt zwei Grundformen.

\textbf{Vektordaten} beschreiben einzelne Objekte durch Koordinaten. Für diese Arbeit
reichen drei Typen:

\begin{itemize}
    \item Ein \emph{Punkt} ist eine einzelne Koordinate, zum Beispiel die GPS-Ortung eines
    besenderten Bären oder der Standort eines Geräts.
    \item Eine \emph{Linie} ist eine Folge von Punkten, zum Beispiel der Bewegungspfad eines
    Bären über einen Tag.
    \item Ein \emph{Polygon} ist eine geschlossene Fläche, zum Beispiel ein
    Untersuchungsgebiet oder eine von Nutzenden eingezeichnete Fläche.
\end{itemize}

\textbf{Rasterdaten} bestehen aus einem Gitter von Zellen mit je einem Wert.
Satellitenbilder gehören dazu, ebenso Höhenmodelle und die Kartenbilder, aus denen
Web-Karten zusammengesetzt sind.

Für die Darstellung ist der Unterschied wichtig. Ein Rasterbild wird gezeichnet, wie es
ist. Vektordaten dagegen enthalten keine Angaben zum Aussehen: Farbe, Größe und Form
werden erst beim Zeichnen festgelegt. Deshalb lassen sich dieselben Daten je nach Frage
unterschiedlich darstellen.

\subsection{Koordinatenbezugssysteme und GeoJSON}
\label{subsec:crs}

Eine Koordinate allein ist mehrdeutig. Erst mit dem zugehörigen Bezugssystem bezeichnet
sie einen bestimmten Ort \cite{GISGrundlagen}. Ein solches Koordinatenreferenzsystem legt
fest, wie die Erdfigur angenähert wird, wo der Ursprung liegt und in welcher Einheit und
Reihenfolge die Achsen gezählt werden. Die gebräuchlichen Systeme stehen im \emph{EPSG
Geodetic Parameter Dataset}, einem Katalog der International Association of Oil \& Gas
Producers. Jedes System hat dort eine eindeutige Nummer \cite{EPSG}. In dieser Arbeit
treffen zwei Systeme aufeinander.

Das \textbf{World Geodetic System 1984} (WGS~84, EPSG:4326) ist ein geographisches
Koordinatenreferenzsystem. Seine Achsen sind die geodätische Breite nach Norden und die
geodätische Länge nach Osten, beide in Grad. In diesem System arbeiten GPS-Empfänger, und
GeoJSON schreibt es vor. Alle Daten, die die Karte erreichen, liegen darin vor.

Die Rohdaten der Karpaten-Studie liegen dagegen in \textbf{EPSG:3844} vor, amtlich
\enquote{Pulkovo 1942(58) / Stereo70}. Das ist die rumänische Landesprojektion. Sie bildet
die Erdoberfläche über eine schiefachsige stereographische Projektion auf eine Ebene ab,
mit Ursprungsbreite 46° und Mittelmeridian 25°, und gibt Positionen in Metern an.

Der Unterschied ist grundsätzlich: WGS~84 beschreibt Positionen auf einer gekrümmten
Fläche in Winkeln, EPSG:3844 bildet sie auf eine Ebene ab. Jede solche Abbildung verzerrt.
Eine Projektion kann nur eine Eigenschaft erhalten -- Winkel, Flächen oder Entfernungen --
aber nie alle zugleich. Die stereographische Projektion ist winkeltreu und eignet sich
deshalb für Gebiete, die in alle Richtungen etwa gleich weit reichen \cite{Snyder1987}.

Daraus folgt, warum umgerechnet werden muss. Eine Landesprojektion ist auf ein Gebiet
zugeschnitten: Innerhalb Rumäniens sind die Verzerrungen gering und Entfernungen direkt in
Metern ablesbar. Außerhalb wachsen sie schnell an, und für eine Weltkarte ist das System
unbrauchbar.

Die Umrechnung erfolgt deshalb schon beim Lesen aus der
Datenbank.\footnote{\texttt{wildlife\_db\_mariadb.py}: Anlage der Transformation im
Konstruktor, Anwendung je Koordinate in \texttt{\_transform\_coordinates()}.} Die
Datenzugriffsschicht legt die Transformation einmal an und wendet sie auf jede gelesene
Koordinate an.

Dabei ist ein Detail zu beachten. Das EPSG-Register führt die Achsen von EPSG:4326 in der
Reihenfolge Breite, dann Länge. GeoJSON schreibt die umgekehrte Reihenfolge vor, wie der
folgende Abschnitt zeigt. Die Transformation wird deshalb mit dem Parameter
\texttt{always\_xy} angelegt. Er erzwingt die von GeoJSON verlangte Reihenfolge Länge vor
Breite.

Dieser Konflikt ist nicht nur theoretisch. Bei einer externen Datenquelle trat er
tatsächlich auf: Die gelieferten Flächen lagen auf der Karte in einem völlig anderen
Erdteil. Abschnitt~\ref{sec:externe-quellen} beschreibt den Fall und seine Behebung.

Dass die Umrechnung zentral in der Datenzugriffsschicht stattfindet, hat eine praktische
Folge: Alle nachgelagerten Komponenten -- Karte, Simulation, Statistik -- arbeiten nur mit
WGS~84 und müssen die Projektion nicht kennen.

\paragraph{GeoJSON als Austauschformat.}
\label{subsec:geojson}

Für den Austausch von Vektordaten zwischen Server und Browser wird in dieser Arbeit
durchgängig GeoJSON verwendet. Die Internet Engineering Task Force (IETF) hat es als
RFC~7946 standardisiert \cite{RFC7946}. Es baut auf JSON (JavaScript Object Notation) auf,
einem textbasierten Format aus verschachtelten Schlüssel-Wert-Paaren.

Der Aufbau ist einfach. Ein \texttt{Feature} besteht aus einer Geometrie und einem
Eigenschaftsobjekt (\texttt{properties}) mit beliebigen Schlüssel-Wert-Paaren. Mehrere
Features bilden zusammen eine \texttt{Feature\-Collection}.
Listing~\ref{lst:geojson-grundlagen} zeigt ein Punkt-Feature.

\begin{lstlisting}[caption={Aufbau eines GeoJSON-Punkt-Features},
                   label={lst:geojson-grundlagen}]
{
  "type": "Feature",
  "properties": { "bear_id": "10", "timestamp": "2013-06-14T08:00:00" },
  "geometry": { "type": "Point", "coordinates": [25.041, 45.663] }
}
\end{lstlisting}

Das Format passt aus zwei Gründen zu dieser Arbeit. Erstens können beide Seiten es direkt
verarbeiten: Python erzeugt es als verschachteltes Dictionary, JavaScript liest es ohne
Umwandlung. Zweitens nimmt das freie \texttt{properties}-Objekt fachliche Angaben direkt am
Objekt auf, etwa die Bären-Kennung oder den Zeitstempel. Die Kartenbibliothek kann die
Darstellung später von diesen Angaben abhängig machen.

Zwei Festlegungen des Standards sind praktisch wichtig. Erstens schreibt RFC~7946 WGS~84
als Bezugssystem vor. Zweitens legt er die Reihenfolge der Koordinaten fest: Die ersten
beiden Elemente sind \enquote{longitude and latitude [\ldots] precisely in that order}
\cite{RFC7946}, also Länge vor Breite.

Das ist umgekehrt zu der Reihenfolge, in der man Koordinaten gewöhnlich nennt, und
widerspricht außerdem dem EPSG-Register. Dienste, die sich an das Register halten, liefern
GeoJSON, dessen Geometrien an einer völlig falschen Stelle der Erde liegen. Der Fehler
fällt nicht an den Daten selbst auf, sondern nur daran, dass nichts an der erwarteten
Stelle erscheint.

\subsection{Kacheln und Höhenmodelle}
\label{subsec:kacheln}

Eine Web-Karte lädt ihren Hintergrund nicht als ein großes Bild, sondern in quadratischen
Ausschnitten, den \emph{Kacheln}. Alle Anbieter verwenden dafür dasselbe Schema.

Auf Zoomstufe~0 zeigt eine einzige Kachel die ganze Welt. Mit jeder weiteren Zoomstufe
wird jede Kachel in vier Teile geteilt, auf Stufe~$z$ gibt es also $4^{z}$ Kacheln. Jede
Kachel hat die Adresse $(z, x, y)$, die direkt in eine URL eingesetzt wird. Der Browser
lädt nur die Kacheln, die gerade sichtbar sind, und holt beim Verschieben die neuen nach.

Das Raster braucht eine Projektion, die sich gleichmäßig in Quadrate teilen lässt. Üblich
ist die Web-Mercator-Projektion (EPSG:3857). Die Anwendung selbst rechnet nicht in diesem
System; die Kartenbibliothek erledigt das intern.

Daraus folgt eine Einschränkung, die leicht übersehen wird: Jeder Dienst hält Kacheln nur
bis zu einer bestimmten Zoomstufe vor. Zoomt man weiter hinein, muss der Browser die
vorhandenen Kacheln vergrößern. Kennt er diese Obergrenze nicht, fordert er Kacheln an, die
es nicht gibt. Kapitel~\ref{chap:implementierung} zeigt einen solchen Fall.

\textbf{Rasterkacheln} sind fertige Bilder, meist PNG oder JPEG. Der Server legt das
Aussehen fest, der Browser zeichnet das Bild nur noch. Satellitenbilder gibt es nur in
dieser Form, weil sie Fotografien sind.

\textbf{Vektorkacheln} enthalten statt eines Bildes die Geometrien selbst, in einem binären
Format nach der Mapbox Vector Tile Specification \cite{MVTSpec}. Hier bestimmt erst der
Browser das Aussehen. Farben, Linienstärken und Beschriftungen lassen sich deshalb ändern,
ohne neue Daten zu laden -- dieselben Kacheln ergeben eine helle und eine dunkle Karte.

\paragraph{Höhenmodelle.}
\label{subsec:hoehenmodelle}

Für die dreidimensionale Darstellung des Geländes braucht die Karte ein digitales
Höhenmodell (Digital Elevation Model, DEM). Es kommt ebenfalls als Rasterkachel, kodiert
aber keine Farbe, sondern eine Höhe: Die drei Farbkanäle eines Bildpunktes werden gemeinsam
als eine Zahl gelesen, die die Geländehöhe in Metern angibt.

Daraus baut der Browser ein Höhengitter und hebt die Kartengeometrie entsprechend an. Aus
denselben Daten lässt sich eine Schummerung berechnen, die das Relief durch simulierten
Lichteinfall auch in der Draufsicht sichtbar macht. Beide Darstellungsformen dieser Arbeit
beruhen damit auf derselben Datengrundlage.

% =====================================================================
\section{Web-Kartenbibliotheken und ihre Einbindung}
\label{sec:bibliotheken}

Die eigentliche Darstellung der Karte im Browser übernimmt eine JavaScript-Bibliothek.
Dieser Abschnitt stellt die beiden Bibliotheken gegenüber, die für diese Arbeit relevant
sind: Leaflet, auf dem das Ausgangssystem beruhte, und MapLibre GL JS, auf das es
umgestellt wurde. Anschließend wird der Anwendungsrahmen behandelt, in den beide
eingebettet werden.

\subsection{Leaflet und Folium}
\label{subsec:leaflet}

Leaflet ist eine verbreitete, quelloffene JavaScript-Bibliothek für interaktive Karten. Sie
verwaltet die Kacheln des Kartenhintergrunds, reagiert auf Zoomen und Verschieben und
stellt Vektorobjekte wie Marker, Linien und Flächen bereit.

Für das Zeichnen dieser Vektorobjekte verwendet Leaflet standardmäßig einen SVG-Renderer;
über die Option \texttt{preferCanvas} lässt sich alternativ ein Canvas-Renderer wählen
\cite{LeafletDocs}. Im SVG-Fall erhält jedes Objekt ein eigenes Element im Document Object
Model des Browsers. Der Browser muss diese Elemente verwalten und bei jeder Kartenbewegung
neu positionieren; der Aufwand wächst somit mit der Zahl der Objekte. Der Canvas-Renderer
vermeidet die Einzelelemente, zeichnet aber ebenfalls auf der CPU.

Leaflet ist eine reine JavaScript-Bibliothek. Damit sie aus einer Python-Anwendung heraus
nutzbar wird, existiert der Wrapper Folium: Er nimmt Python-Objekte entgegen und erzeugt
daraus eine vollständige HTML-Seite mit eingebettetem Leaflet-Code.

Die Einbindung dieser Seite in eine Streamlit-Anwendung erfolgt über das Paket
\texttt{streamlit-folium}. Es bietet zwei Funktionen an, deren Unterschied für das
Verständnis des Ausgangssystems wesentlich ist \cite{StreamlitFolium}:

\begin{itemize}
    \item \texttt{st\_folium()} wird als \enquote{bi-directional Component} beschrieben.
    Die Funktion liefert ein Dictionary zurück, über das Benutzerinteraktionen - etwa
    gezeichnete Objekte oder angeklickte Elemente - an Python zurückgemeldet werden.
    \item \texttt{folium\_static()} stellt die Karte über deren HTML-Repräsentation dar und
    ist laut Dokumentation eine echte Teilmenge der Funktionalität von
    \texttt{st\_folium()}. Es gibt keinen Rückgabewert; Benutzereingaben erreichen die
    Anwendung nicht.
\end{itemize}

Eine mit \texttt{folium\_static()} eingebundene Karte kann folglich nichts entgegennehmen,
was der Nutzer auf ihr tut. Dieser Unterschied erklärt die in
Kapitel~\ref{chap:einleitung} beschriebene Ungleichheit der beiden Ausgangskarten und wird
in Kapitel~\ref{chap:analyse} anhand des Quellcodes belegt.

\subsection{MapLibre GL JS und Begründung der Technologiewahl}
\label{subsec:maplibre}

Mapbox GL JS geht einen anderen Weg. Die Bibliothek verwaltet Vektorobjekte nicht als
Elemente im Dokumentbaum, sondern übergibt sie als Geometriedaten an die Grafikkarte und
lässt sie dort über die Browser-Schnittstelle WebGL zeichnen. Damit wurden Vektorkacheln
im Browser praktisch nutzbar.

Im Dezember 2020 stellte Mapbox die Bibliothek auf eine nicht quelloffene Lizenz um. Aus
dem letzten frei lizenzierten Stand entstand MapLibre GL JS als quelloffener Fork, zunächst
als direkter Ersatz gedacht. MapLibre GL JS steht unter der 3-Clause-BSD-Lizenz
\cite{MapLibreRepo}. Die eigene Dokumentation beschreibt es als \enquote{a TypeScript
library that uses WebGL to render interactive maps from vector tiles in a browser}
\cite{MapLibreDocs}.

Die Lizenzfrage ist für diese Arbeit wichtig. Eine quelloffene Bibliothek ohne Bindung an
einen kommerziellen Anbieter braucht keinen Zugangsschlüssel, verursacht keine laufenden
Kosten und hängt nicht von fremden Geschäftsbedingungen ab. Für ein Forschungsprojekt mit
ungewisser Laufzeit ist das eine Voraussetzung.

\paragraph{Datenquellen und Ebenen.}

MapLibre GL JS trennt Daten und Darstellung. Das ist sein zentrales Gestaltungsprinzip.

Eine \emph{Datenquelle} (\texttt{source}) verweist auf einen Datenbestand, etwa eine
GeoJSON-FeatureCollection oder einen Vektorkachel-Dienst. Sie legt fest, \emph{was}
dargestellt wird, nicht wie es aussieht.

Eine \emph{Ebene} (\texttt{layer}) greift auf eine Datenquelle zu und legt fest, \emph{wie}
deren Inhalt gezeichnet wird: als Kreis, als Linie oder als gefüllte Fläche, jeweils mit
Farbe und Größe. Mehrere Ebenen können dieselbe Datenquelle unterschiedlich darstellen.

Wichtig ist, dass das Aussehen von den Eigenschaften der Daten abhängen darf. Eine Ebene
kann die Farbe eines Punktes aus dessen \texttt{properties} ableiten, statt sie fest
vorzugeben. Eine Unterscheidung wie die zwischen gemessenen und simulierten Punkten wird so
einmal in den Daten hinterlegt und von der Karte automatisch umgesetzt.

Die gesamte Kartenkonfiguration steht in einem Stildokument, dessen Aufbau die MapLibre
Style Spec festlegt \cite{MapLibreStyleSpec}. Eine Änderung daran wirkt sofort, ohne dass
Daten neu geladen werden.

MapLibre GL JS unterstützt außerdem Höhenmodelle als Datenquelle und eine neigbare Kamera.
Damit ist die dreidimensionale Darstellung des Geländes möglich \cite{MapLibreDocs}.

\paragraph{Gegenüberstellung der Bibliotheken.}
\label{subsec:vergleich}

Tabelle~\ref{tab:bibliotheken} stellt die wichtigsten Unterschiede gegenüber. Verglichen
werden Funktionsweise und vorhandene Funktionen; wie sich die Darstellung in der Praxis
schlägt, bewertet Kapitel~\ref{chap:diskussion}.

\begin{table}[htbp]
\centering
\caption{Gegenüberstellung der betrachteten Kartenbibliotheken}
\label{tab:bibliotheken}
\small
\begin{tabular}{p{3.8cm} p{4.6cm} p{4.6cm}}
\toprule
\textbf{Kriterium} & \textbf{Leaflet (über Folium)} & \textbf{MapLibre GL JS} \\
\midrule
Zeichnen von Vektorobjekten
  & SVG (Standard) oder Canvas, jeweils auf der CPU
  & WebGL, auf der Grafikkarte \\
\addlinespace
Verwaltung einzelner Objekte
  & Bei SVG ein Element je Objekt im Dokumentbaum
  & Keine Einzelelemente im Dokumentbaum \\
\addlinespace
Kachelarten
  & Rasterkacheln
  & Raster- und Vektorkacheln \\
\addlinespace
Änderung des Aussehens zur Laufzeit
  & Neuaufbau der betroffenen Objekte
  & Eigenschaft im Stildokument ändern \\
\addlinespace
Dreidimensionale Geländedarstellung
  & Nicht vorgesehen
  & Über Höhenmodell als Datenquelle \\
\addlinespace
Einbindung in Python
  & Folium als ausgereifter Wrapper
  & Kein vergleichbarer Wrapper \\
\bottomrule
\end{tabular}
\end{table}

Die Wahl fiel aus drei Gründen auf MapLibre GL JS. Erstens entfallen die Einzelelemente im
Dokumentbaum, deren Verwaltung mit der Objektzahl aufwendiger wird -- genau das wird zum
Problem, wenn alle Datenbestände auf einer Karte liegen. Zweitens lassen sich dieselben
Daten dank der Trennung von Datenquelle und Ebene ohne Neuladen unterschiedlich darstellen
und einzelne Ebenen ein- und ausblenden. Drittens gibt es Satellitengrundkarten und
dreidimensionale Geländedarstellung, die Leaflet nicht bietet.

Dem steht der Nachteil aus der letzten Tabellenzeile gegenüber: Für MapLibre GL JS gibt es
keinen Python-Wrapper wie Folium. Die Einbindung in die Streamlit-Anwendung muss deshalb
selbst entwickelt werden. Welche Möglichkeiten es dafür gibt, zeigt der folgende Abschnitt;
welche gewählt wurde und warum, steht in Kapitel~\ref{chap:implementierung}.

\subsection{Streamlit als Anwendungsrahmen}
\label{subsec:streamlit}

Die Anwendung dieser Arbeit ist in Python geschrieben und nutzt Streamlit als Rahmen für
die Benutzeroberfläche. Streamlit erzeugt aus einem gewöhnlichen Python-Skript eine
Webanwendung: Aufrufe wie \texttt{st.selectbox()} oder \texttt{st.button()} erzeugen
Bedienelemente, deren Werte das Skript unmittelbar weiterverarbeiten kann.

Drei Eigenschaften dieses Rahmens bestimmen den Aufbau der späteren Implementierung und
werden deshalb hier eingeführt.

\paragraph{Das Ausführungsmodell.}

Streamlit kennt keine Ereignisbehandlung im klassischen Sinn. Stattdessen gilt:
\enquote{Streamlit runs your script from top to bottom at every user interaction or code
change} \cite{StreamlitCaching}.

Jede Interaktion - ein Klick, die Auswahl in einer Liste, das Verschieben eines Reglers --
führt also dazu, dass das gesamte Skript erneut von der ersten bis zur letzten Zeile
abläuft. Die Bedienelemente geben dabei ihre aktuellen Werte zurück, und die Oberfläche
wird aus dem Ergebnis neu aufgebaut.

Dieses Modell macht den Programmablauf einfach nachvollziehbar, hat aber eine unmittelbare
Folge: Jede aufwendige Berechnung im Skript wird bei jeder Interaktion wiederholt - auch
dann, wenn sie mit der ausgelösten Änderung nichts zu tun hat. Eine Datenbankabfrage für
die eine Datenebene läuft also auch dann erneut, wenn der Nutzer einen Filter einer ganz
anderen Datenebene verstellt.

\paragraph{Zwischenspeicherung.}

Gegen dieses Problem stellt Streamlit zwei Dekoratoren bereit, die sich in ihrem
Einsatzzweck unterscheiden \cite{StreamlitCaching}:

\begin{itemize}
    \item \texttt{st.cache\_data} ist für Daten gedacht, also für serialisierbare
    Rückgabewerte wie Tabellen, Zahlenfelder oder Zeichenketten. Jeder Aufruf erhält eine
    eigene Kopie des Ergebnisses, wodurch ungewollte Veränderungen durch einen Aufrufer
    andere Aufrufer nicht beeinflussen.
    \item \texttt{st.cache\_resource} ist für globale Betriebsmittel gedacht, die sich
    nicht serialisieren lassen - etwa Datenbankverbindungen oder geladene Modelle. Alle
    Durchläufe und alle Sitzungen teilen sich dasselbe Objekt, ohne es zu kopieren. Das
    Objekt muss daher nebenläufigkeitssicher sein, und Änderungen daran wirken sich auf
    alle Nutzer aus.
\end{itemize}

Die Unterscheidung ist nicht formal: Ein geladenes Simulationsmodell gehört in die zweite
Kategorie, ein daraus berechnetes Ergebnis in die erste. Kapitel~\ref{chap:implementierung}
zeigt, wo in der Anwendung welcher Dekorator eingesetzt wird.

\paragraph{Einbindung eigener Komponenten.}
\label{par:komponenten}

Streamlit bringt kein Kartenwidget mit dem benötigten Funktionsumfang mit. MapLibre GL JS
muss deshalb als eigene Komponente eingebunden werden. Dafür bestehen zwei Wege, die sich
grundlegend unterscheiden \cite{StreamlitComponents}.

Gemeinsam ist beiden, dass die Komponente als eigenständige Webseite in einem
\texttt{iframe} dargestellt wird. Sie läuft damit in einem eigenen Dokument und hat keinen
direkten Zugriff auf die umgebende Streamlit-Seite.

\textbf{Statische Komponenten} werden über \texttt{components.html()} erzeugt. Die Funktion
nimmt eine Zeichenkette mit HTML entgegen und bettet sie in einen \texttt{iframe} ein. Die
Übertragung verläuft nur in eine Richtung: Die Komponente zeigt an, was Python ihr
mitgibt, meldet aber nichts zurück.

\textbf{Bidirektionale Komponenten} werden über \texttt{declare\_component()} erzeugt. Beim
Aufruf übergebene Argumente erreichen die Komponente; ruft deren JavaScript-Code
\texttt{Streamlit.setComponentValue()} auf, gelangt der übergebene Wert zurück nach
Python. Streamlit führt daraufhin das Skript erneut aus, und der Aufruf der Komponente
liefert in diesem Durchlauf den neuen Wert. Da nur ein einzelner Wert übergeben werden
kann, müssen mehrere Angaben in einem Objekt zusammengefasst werden.

Der Unterschied entspricht strukturell dem zwischen \texttt{folium\_static()} und
\texttt{st\_folium()} aus Abschnitt~\ref{subsec:leaflet}: Auch dort steht eine
Einbahnstraße einer Verbindung mit Rückkanal gegenüber. Welcher Weg in dieser Arbeit
gewählt wurde, welche Folgen das hat und wie der fehlende Rückkanal umgangen wird,
behandelt Kapitel~\ref{chap:implementierung}.

% =====================================================================
\section{Bewegungsmodellierung und Kamerafallendaten}
\label{sec:fachliche-grundlagen}

Dieser Abschnitt führt die fachlichen Grundlagen ein, die für die Erweiterung des
Bewegungsmodells und für den statistischen Vergleich gebraucht werden.

\subsection{Bewegungsökologie und Step-Selection-Funktionen}
\label{subsec:bewegungsoekologie}

\paragraph{Trajektorie, Schritt und Drehwinkel.}

In der Bewegungsökologie wird die Bewegung eines Tieres als Folge einzelner Schritte
beschrieben. Ein \emph{Schritt} ist der Vektor zwischen zwei aufeinanderfolgenden
Positionen $x_{t-1}$ und $x_{t}$ im Zeitabstand $\Delta t$. Zwei Größen beschreiben ihn
vollständig:

\begin{itemize}
    \item die \emph{Schrittlänge} $l$, also die zurückgelegte Entfernung,
    \item der \emph{Drehwinkel} $\theta$, also die Richtungsänderung gegenüber dem
    vorherigen Schritt.
\end{itemize}

Fast alle Bewegungsmodelle setzen hier an. Die Zerlegung verlangt allerdings gleichmäßige
Zeitabstände zwischen den Positionen, und das ist bei echten Telemetriedaten wegen
Messlücken nicht durchgängig gegeben.

\paragraph{Korrelierter Zufallspfad.}

Das einfachste Modell ist der einfache Zufallspfad: Jeder Schritt bekommt eine neue,
gleichverteilt gezogene Richtung. Echte Tiere laufen nicht so. Sie behalten eine Richtung
über mehrere Schritte bei. Diese Eigenschaft heißt Richtungstreue.

Der \emph{korrelierte Zufallspfad} bildet sie ab: Der Drehwinkel wird nicht
gleichverteilt gezogen, sondern aus einer Verteilung, die sich um null häuft. Üblich ist
dafür die von-Mises-Verteilung, das Gegenstück zur Normalverteilung für Winkel. Ihre Dichte
lautet

\begin{equation}
    p(\theta) \;=\; \frac{1}{2\pi I_{0}(\kappa)}\,
    \exp\!\bigl(\kappa \cos(\theta - \mu)\bigr),
    \label{eq:vonmises}
\end{equation}

mit der mittleren Richtung $\mu$, dem Konzentrationsparameter $\kappa$ und der
modifizierten Besselfunktion erster Art nullter Ordnung $I_{0}$.

$\kappa$ steuert das Verhalten vollständig. Bei $\kappa = 0$ ist die Verteilung über den
Kreis gleichmäßig, das Tier läuft also richtungslos umher. Je größer $\kappa$, desto stärker
konzentriert sich die Verteilung auf die bisherige Richtung und desto gestreckter wird die
Bewegung.

Wichtig für die Anwendung: $\kappa$ lässt sich aus Daten schätzen. Für beobachtete
Drehwinkel $\theta_{1},\dots,\theta_{n}$ ist die mittlere resultierende Länge

\begin{equation}
    \rho \;=\; \frac{1}{n}\sum_{j=1}^{n} \cos\theta_{j}.
    \label{eq:rho}
\end{equation}

Zwischen $\rho$ und $\kappa$ gilt $A(\kappa) = I_{1}(\kappa)/I_{0}(\kappa) = \rho$; gesucht
ist also $\kappa = A^{-1}(\rho)$. Diese Umkehrung lässt sich nicht geschlossen angeben und
wird in der Praxis genähert.

Anschaulich misst $\rho$, wie geradlinig die Bewegung ist: $+1$ heißt, jeder Schritt setzt
die Richtung des vorherigen genau fort, $0$ heißt richtungslos, negative Werte heißen
überwiegend Hin und Her.

\paragraph{Step-Selection-Funktionen.}

Step-Selection-Funktionen (SSF) verbinden Bewegung mit Umweltpräferenz. Das Prinzip ist
einfach: Jedem beobachteten Schritt werden $k$ \emph{mögliche} Schritte gegenübergestellt.
Sie starten am selben Punkt, führen aber zu zufällig gezogenen Zielen. Aus dem Vergleich
der Umwelteigenschaften am echten Ziel mit denen an den Alternativen lässt sich schätzen,
was bevorzugt und was gemieden wird \cite{Thurfjell2014}.

Die Methode ist in der Bewegungsökologie etabliert und wird unter anderem für Habitatwahl,
Wanderkorridore und Mensch-Wildtier-Konflikte eingesetzt \cite{Thurfjell2014}. Wichtig ist
dabei, wie die Alternativschritte gezogen werden: Das Ziehverfahren beeinflusst die
Schätzung systematisch \cite{Forester2009}.

Eine Weiterentwicklung ist die integrierte Step-Selection-Analyse. Sie gibt die Annahme auf,
dass Bewegungseigenschaften und Umweltpräferenz unabhängig voneinander sind, und schätzt
beides gemeinsam \cite{Avgar2016}. Das R-Paket \texttt{amt} stellt dafür Werkzeuge bereit
\cite{Signer2019}.

\paragraph{Einordnung für diese Arbeit.}

Das SSF-Prinzip wird hier in einer bestimmten Richtung verwendet, und das ist für die
Abgrenzung wichtig.

Die klassische Anwendung ist \emph{analytisch}: Aus beobachteten Schritten und erzeugten
Alternativen werden Präferenzparameter geschätzt, meist über bedingte logistische
Regression. Dieser Schritt findet hier nicht statt. Es wird keine Regression gerechnet und
es werden keine Präferenzkoeffizienten geschätzt.

Verwendet wird nur der \emph{erzeugende} Teil: An jeder Position gibt es eine Menge
erreichbarer Kandidatenschritte, aus der einer ausgewählt wird. Die Auswahl erfolgt
gewichtet, nach Richtungstreue und in der erweiterten Fassung zusätzlich nach dem Gelände.
Das Ergebnis ist kein Schätzwert, sondern ein simulierter Pfad.

Bei der Bewertung der Ergebnisse ist das zu beachten: Die erzeugten Pfade sind keine aus
Daten abgeleiteten Aussagen über Präferenzen, sondern Realisierungen eines Modells, dessen
Parameter teils aus den Karpaten-Daten gemessen und teils gesetzt sind.

\subsection{Gewichtete Schwerpunktbildung}
\label{subsec:schwerpunkt}

Ein letzter Baustein stammt nicht aus der Bewegungsökologie, sondern aus einem anderen
Anwendungsgebiet.

Geyer, Melzer, Dünnweber und Sarkar beschreiben ein Verfahren zur Lokalisierung von
Fußpositionen auf Bodensensormatten \cite{GeyerMelzerDuennweber2026}. Die Sensorfläche ist
dort in Felder unterteilt, die entlang ihrer Diagonalen nochmals in Dreiecke zerlegt sind.
Betritt ein Fuß die Fläche, werden mehrere dieser Dreiecke unterschiedlich stark aktiviert.

Die Position wird dann nicht als Mittelpunkt der aktivierten Fläche bestimmt, sondern als
gewichteter Schwerpunkt \cite[Formel~(1)]{GeyerMelzerDuennweber2026}:

\begin{equation}
    P \;=\; \frac{\sum_{i=1}^{n} w_{i}\, p_{i}}{\sum_{i=1}^{n} w_{i}} .
    \label{eq:schwerpunkt}
\end{equation}

Dabei ist $P$ die gesuchte Position, $p_{i}$ der Mittelpunkt des $i$-ten aktivierten
Dreiecks und $w_{i}$ dessen Aktivierungswert. Die Gewichtung bewirkt, dass die berechnete
Position sich zu den stärker aktivierten Bereichen hin verschiebt, also dorthin, wo der
größere Teil des Fußes aufliegt.

Für die vorliegende Arbeit ist weniger der konkrete Anwendungsfall von Bedeutung als die
zugrunde liegende Idee: Aus einer Menge unterschiedlich stark bewerteter Flächenstücke wird
ein einzelner, nach diesen Bewertungen gewichteter Bezugspunkt bestimmt. Dieselbe Aufgabe
stellt sich bei der Richtungswahl in strukturiertem Gelände, wenn aus einer Menge
unterschiedlich attraktiver Geländezellen ein Ziel abgeleitet werden soll.
Kapitel~\ref{chap:implementierung} beschreibt die Übertragung.

\subsection{Kamerafallen und Aktivitätsmuster}
\label{subsec:kamerafallen}

\paragraph{Detektionen statt Trajektorien.}

Kamerafallen liefern eine ganz andere Datenform als GPS-Telemetrie. Eine Kamera löst aus,
wenn sich vor ihr etwas bewegt, und speichert Bild und Zeitpunkt. Daraus folgen drei
Einschränkungen.

Erstens lässt sich das fotografierte Tier meist nicht als Individuum bestimmen. Zwei
Aufnahmen können dasselbe Tier zeigen oder zwei verschiedene.

Zweitens ist zwischen zwei Aufnahmen kein Bewegungszusammenhang bekannt. Aus Detektionen an
verschiedenen Kameras lässt sich kein Weg rekonstruieren, weil weder feststeht, ob es
dasselbe Tier war, noch welchen Weg es genommen hätte.

Drittens wird nur punktuell beobachtet. Erfasst wird, was sich im Bereich der Kamera bewegt;
über alles dazwischen gibt es keine Information.

Deshalb gibt es in dieser Arbeit keinen räumlichen Pfadvergleich zwischen Simulation und
japanischen Realdaten: Auf der Seite der Realdaten gibt es keine Pfade, mit denen man
vergleichen könnte. Der Vergleich wird stattdessen auf die Zeit verlagert.

\paragraph{Aktivitätsmuster als zirkuläre Daten.}

Ohne Individualerkennung lässt sich aus Kamerafallendaten trotzdem das zeitliche
Aktivitätsmuster ableiten. Die Idee: Die Uhrzeiten der Auslösungen sind eine Stichprobe aus
der Tagesaktivität der Population \cite{Rowcliffe2014}.

Uhrzeiten sind dabei keine gewöhnlichen Zahlen, sondern \emph{zirkuläre} Daten. Der Tag ist
ein geschlossener Kreis: 23:59~Uhr und 00:01~Uhr liegen zwei Minuten auseinander, als Zahlen
aber fast 24~Stunden. Wer das nicht berücksichtigt, bekommt genau an der Tagesgrenze falsche
Ergebnisse -- besonders bei nachtaktiven Arten, bei denen sich viele Beobachtungen um
Mitternacht häufen.

Die Uhrzeit wird deshalb in einen Winkel umgerechnet, wobei 24~Stunden einem vollen Umlauf
von $2\pi$ entsprechen. Aus diesen Winkeln wird dann eine Dichtefunktion geschätzt, die die
Aktivität über den Tag beschreibt. Dafür dient eine Kerndichteschätzung mit einem
zirkulären Kern \cite{RidoutLinkie2009}.

\paragraph{Der Überlappungskoeffizient $\Delta$.}

Liegen zwei solche Dichtefunktionen vor -- etwa für zwei Arten oder, wie hier, für Realdaten
und Simulation --, stellt sich die Frage, wie ähnlich sie sind. Ridout und Linkie schlagen
dafür den Überlappungskoeffizienten $\Delta$ vor \cite{RidoutLinkie2009}. Er ist die Fläche
unter dem Minimum beider Dichten:

\begin{equation}
    \Delta \;=\; \int_{0}^{24} \min\bigl(f(t),\, g(t)\bigr)\, \mathrm{d}t .
    \label{eq:delta}
\end{equation}

Anschaulich ist das die Fläche, die sich beide Kurven teilen. Der Wert liegt zwischen 0
und 1: $\Delta = 0$ heißt, die Aktivitätszeiten überschneiden sich nicht, $\Delta = 1$ heißt
gleicher Tagesrhythmus.

Ridout und Linkie geben mehrere Schätzer für $\Delta$ an und prüfen sie in einer
Simulationsstudie; welcher passt, hängt vom Stichprobenumfang ab \cite{RidoutLinkie2009}.
Die Referenzimplementierung ist das R-Paket \texttt{overlap} \cite{Meredith2016}.

Das Maß ist der Standard für den Vergleich von Tagesrhythmen. Wichtig ist aber, was es
\emph{nicht} sagt: $\Delta$ misst nur die zeitliche Übereinstimmung. Über das räumliche
Verhalten, über zurückgelegte Strecken oder darüber, ob sich zwei Tiere wirklich begegnen,
sagt der Wert nichts.

\paragraph{Stichprobenumfang und Unsicherheit.}

$\Delta$ wird aus einer begrenzten Zahl von Beobachtungen geschätzt und ist deshalb selbst
unsicher. Zwei Punkte sind dabei wichtig.

Zum einen braucht es eine Mindestzahl an Beobachtungen. Eine Untersuchung zu
Aktivitätskurven zeigt, dass sich Mittelwert und Streuung der Schätzung mit wachsender
Stichprobe schnell stabilisieren, und nennt rund 100~Detektionen als Richtwert
\cite{Cove2018}.

Zum anderen sollte man einen Einzelwert nie ohne Unsicherheitsangabe berichten. Üblich ist
das Bootstrap-Verfahren: Aus den vorhandenen Beobachtungen wird wiederholt eine neue
Stichprobe gleicher Größe mit Zurücklegen gezogen, für jede wird $\Delta$ neu berechnet, und
aus der Verteilung dieser Werte ergibt sich ein Konfidenzintervall. Das Verfahren braucht
keine Annahme über die Verteilung und passt deshalb auch zu zirkulären Daten.

Beide Punkte kommen in Kapitel~\ref{chap:implementierung} bei der Berechnung und in
Kapitel~\ref{chap:diskussion} bei der Einordnung wieder vor.

% ====================================================================
\chapter{Analyse der bestehenden Systeme und Ableitung der Anforderungen}
\label{chap:analyse}

Dieses Kapitel beschreibt den Zustand der Systeme zu Beginn dieser Arbeit. Es behandelt die vorhandenen Komponenten, die Datenbestände sowie die übernommenen Modelle. Auf der Grundlage der identifizierten Einschränkungen werden anschließend die Anforderungen an das zu entwickelnde System abgeleitet.

% =====================================================================
\section{Ausgangszustand des Systems}
\label{sec:ausgangszustand}

\subsection{Überblick über die Teilsysteme}
\label{subsec:systemuebersicht}

Zu Beginn dieser Arbeit lagen die Ergebnisse von drei Vorarbeiten vor, die unabhängig
voneinander entstanden sind. Nur eine davon war eine Anwendung mit Benutzeroberfläche.
Tabelle~\ref{tab:teilsysteme} nennt die drei Arbeiten,
Abbildung~\ref{fig:ausgangszustand} zeigt, auf welche Daten sie zugreifen.

\begin{table}[htbp]
\centering
\caption{Übersicht der vorgefundenen Vorarbeiten}
\label{tab:teilsysteme}
\small
\begin{tabular}{L{2.9cm} L{5.2cm} L{4.6cm}}
\toprule
\textbf{Vorarbeit} & \textbf{Gegenstand} & \textbf{Datenhaltung} \\
\midrule
Marcel Lorenz \cite{LorenzBA}
  & Kartenanwendung mit zwei Seiten: Gerätestandorte in Japan sowie
    Bewegungspfade der Karpatenbären
  & MariaDB; Bilder als Dateien im Dateisystem, in der Datenbank nur deren Pfade \\
\addlinespace
Fabian Bräu \cite{BreuBA}
  & Modell zur Erzeugung alternativer Bewegungsschritte
  & eigene SQLite-Datei und CSV-Datei \\
\addlinespace
Anthony Lelchat \cite{LelchatBA}
  & Übernahme neuer Kameraaufnahmen aus einem E-Mail-Postfach in die Datenbank
  & dieselbe MariaDB wie die Kartenanwendung \\
\bottomrule
\end{tabular}
\end{table}

\begin{figure}[htbp]
\centering
\begin{tikzpicture}[
    font=\footnotesize,
    >={Stealth[length=2.2mm,width=1.6mm]},
    every node/.style={align=center, inner sep=2pt},
    box/.style={draw=black!80, fill=white, line width=0.5pt,
                minimum height=10mm, text width=#1},
    box/.default=34mm,
    db/.style={cylinder, shape border rotate=90, aspect=0.16, draw=black!80,
               fill=black!4, line width=0.5pt, minimum width=#1,
               minimum height=12mm, text width=#1},
    db/.default=24mm,
    flow/.style={->, line width=0.6pt, draw=black!85},
    ref/.style={->, densely dashed, line width=0.5pt, draw=black!60},
    lbl/.style={font=\scriptsize, text=black!75, inner sep=1.5pt, fill=white},
    head/.style={font=\small\bfseries, text width=40mm},
]

\def\xA{3.1}
\def\xB{8.3}
\def\xC{12.9}
\def\xDB{5.7}

% --- Spaltenköpfe ---
\node[head] (hA) at (\xA,0)
  {Kartenanwendung\\[-1pt]\normalfont\scriptsize\color{black!60}Vorarbeit Lorenz \cite{LorenzBA}};
\node[head] (hB) at (\xB,0)
  {Bildextraktion\\[-1pt]\normalfont\scriptsize\color{black!60}Vorarbeit Lelchat \cite{LelchatBA}};
\node[head] (hC) at (\xC,0)
  {Bewegungsmodell\\[-1pt]\normalfont\scriptsize\color{black!60}Vorarbeit Bräu \cite{BreuBA}};

% --- Reihe 1 ---
\node[box=40mm, minimum height=12mm] (app) at (\xA,-1.75)
  {\textbf{Streamlit-Anwendung}\\[-1pt]\scriptsize zwei getrennte Kartenseiten:\\[-2pt]\scriptsize Japan und Karpaten};
\node[box=34mm, rounded corners=4pt, fill=black!4] (mail) at (\xB,-1.75)
  {E-Mail-Postfach\\[-1pt]\scriptsize Microsoft Graph API};

% --- Reihe 2 ---
\node[box=40mm] (wdb) at (\xA,-3.85)
  {\ttfamily WildlifeDatabase\\[-1pt]\normalfont\scriptsize Zugriff auf die MariaDB};
\node[box=37mm] (ext) at (\xB,-3.85)
  {\scriptsize\ttfamily WildlifePicture-\\[-2pt]\scriptsize\ttfamily Extractor\\[-1pt]\normalfont\scriptsize Texterkennung, Dauerprozess};
\node[box=37mm] (ssm) at (\xC,-3.85)
  {\scriptsize\ttfamily StepSelectionModel\\[-1pt]\normalfont\scriptsize nur Kandidatenschritte};

% --- Reihe 3 ---
\node[db=30mm] (maria) at (\xDB,-5.95)
  {\textbf{MariaDB}\\[-1pt]\scriptsize zentrale Instanz,\\[-2pt]\scriptsize 11 Tabellen};
\node[db=22mm] (fs)     at (1.15,-5.95)      {Dateisystem\\[-1pt]\scriptsize Bilddateien};
\node[db=22mm] (sqlite) at (\xC-1.15,-5.95)  {SQLite\\[-1pt]\scriptsize\ttfamily database.db};
\node[db=22mm] (csv)    at (\xC+1.15,-5.95)  {CSV\\[-1pt]\scriptsize\ttfamily 1\_bears\_RO};

% --- Verbindungen ---
\draw[flow, <->] (app) -- (wdb) node[lbl, midway] {lesen/schreiben};
\draw[flow, <->] (wdb.south) -- ([xshift=-5mm]maria.north);
\draw[flow] (mail) -- (ext);
\draw[flow] (ext.south) -- ([xshift=5mm]maria.north)
  node[lbl, pos=0.55, anchor=north] {schreibt direkt};
\draw[ref] (maria.west) -- (fs.east) node[lbl, midway] {Pfade};
\draw[flow] (ssm.south -| sqlite.north) -- (sqlite.north);
\draw[flow] (ssm.south -| csv.north) -- (csv.north);

% --- isoliertes Bewegungsmodell ---
\begin{scope}[on background layer]
  \node[draw=black!55, dash pattern=on 3pt off 2pt, line width=0.6pt,
        rounded corners=5pt, inner sep=7pt,
        fit=(hC)(ssm)(sqlite)(csv)] (isoC) {};
\end{scope}
\node[font=\scriptsize\itshape, text=black!60, anchor=north]
  at ($(isoC.south)+(0,-2pt)$) {keine Verbindung zu den\\[-2pt]übrigen Teilsystemen};

% --- Legende ---
\begin{scope}[shift={(3.0,-8.1)}, font=\scriptsize]
  \draw[flow] (0,0) -- (0.7,0);
  \node[anchor=west] at (0.8,0) {Datenzugriff};
  \draw[ref] (3.2,0) -- (3.9,0);
  \node[anchor=west] at (4.0,0) {Verweis über Dateipfad};
\end{scope}
\end{tikzpicture}
\caption[Ausgangszustand des Systems]{Ausgangszustand des Systems. Nur die
  Kartenanwendung besitzt eine Oberfläche und greift über
  \texttt{WildlifeDatabase} auf die MariaDB zu.}
\label{fig:ausgangszustand}
\end{figure}

\paragraph{Die vorgefundene Anwendung.} Die Kartenanwendung aus \cite{LorenzBA} hatte genau
zwei Seiten. Über ein Auswahlelement in der Seitenleiste ließ sich zwischen ihnen wechseln.
Jede Seite öffnete eine eigene Leaflet-Karte: entweder mit den Gerätestandorten in Japan
oder mit den Bewegungspfaden der Karpatenbären. Eine gemeinsame Darstellung beider
Datenbestände gab es nicht. Weitere Funktionen enthielt die Anwendung nicht; insbesondere
gehörte die Bewegungssimulation nicht dazu.

Das Bewegungsmodell aus \cite{BreuBA} lag
nur als Quelltext vor und war mit der Kartenanwendung nicht verbunden. Seine Ergebnisse
waren auf der Karte deshalb nicht zu sehen. Angewendet wurde es allein auf die
Karpatendaten, denn diese bestehen aus aufeinanderfolgenden Ortungen eines Tieres und lassen
sich in Schritte zerlegen. Die japanischen Daten geben dagegen Kamerastandorte und
Auslösezeitpunkte wieder und enthalten keine Bewegungsabfolge eines einzelnen Tieres
(Abschnitt~\ref{subsec:japandaten}).

Die Bildextraktion aus \cite{LelchatBA} war ebenfalls keine Anwendung, sondern ein
fortlaufender Prozess: Er holte neue Kameraaufnahmen aus einem E-Mail-Postfach und schrieb
sie in die Datenbank. Eine Kartendarstellung gehörte nicht zu seinem Umfang. Mit der
Kartenanwendung teilte er allein den Datenbestand.

\paragraph{Die einzige systemübergreifende Klammer.} Als Verbindung zwischen den Teilen
diente allein die Klasse \texttt{WildlifeDatabase} (Abbildung~\ref{fig:ausgangszustand}).
Sie kapselt die Verbindung zur zentralen MariaDB und stellt allgemeine Methoden für
Abfragen bereit.

\paragraph{Daraus folgende Aufgabenstellung.} Aus diesem Zustand ergibt sich die zentrale
Aufgabe dieser Arbeit: Die beiden Datenbestände der Kartenanwendung und das bislang
getrennte Bewegungsmodell sollen in einem gemeinsamen System zusammengeführt werden, sodass
die Ergebnisse des Modells auf derselben Karte sichtbar werden wie die realen Daten.



\subsection{Zustand der Kartendarstellung}
\label{subsec:ausgangszustand-karte}

Beide Kartenansichten beruhten auf Leaflet, eingebunden über Folium. Die Art der Einbindung
richtete sich jeweils nach dem Zweck der Seite.

Die japanische Ansicht wurde über \texttt{st\_folium()} eingebunden und nutzte die
Folium-Erweiterung \texttt{Draw}. Damit ließen sich Marker, Rechtecke und Polygone auf der
Karte zeichnen. Die gezeichneten Objekte gelangten als GeoJSON zurück nach Python und wurden
dort auf Knopfdruck in die Datenbank geschrieben.

Die Karpaten-Ansicht diente der Darstellung der Bewegungsdaten und wurde über
\texttt{folium\_static()} eingebunden.\footnote{Zwei Aufrufe in
\texttt{carpathian\_bears.py}, jeweils mit fester Größe von 800$\times$600 Bildpunkten.}
Diese Funktion gibt die Karte als statisches HTML aus; sie besitzt, wie in
Abschnitt~\ref{subsec:leaflet} dargelegt, keinen Rückgabewert.

Jede der beiden Seiten erzeugte eine eigene Karteninstanz mit eigener Konfiguration. Beide
verwendeten den Standard-Renderer von Leaflet, also die Darstellung als Einzelelemente im
Dokumentbaum.

\begin{figure}[htbp]
\centering
\begin{subfigure}{0.48\linewidth}
  \includegraphics[width=\linewidth]{ausgangs-japan.png}
  \caption{Japan-Seite}
\end{subfigure}
\hfill
\begin{subfigure}{0.48\linewidth}
  \includegraphics[width=\linewidth]{ausgangs-karpaten.png}
  \caption{Karpaten-Seite}
\end{subfigure}
\caption[Kartenansichten im Ausgangszustand]{Die beiden Kartenansichten im
  Ausgangszustand. Nur die Japan-Seite nimmt Eingaben entgegen; die
  Werkzeugleiste am linken Kartenrand fehlt in der Karpaten-Ansicht.}
\label{fig:ausgangszustand-karten}
\end{figure}


% =====================================================================
\section{Datenbestände und Datenhaltung}
\label{sec:datenbestaende}

\subsection{GPS-Trajektorien der Karpaten-Studie}
\label{subsec:karpatendaten}

Der umfangreichste Datenbestand stammt aus der in Abschnitt~\ref{sec:motivation}
erwähnten Telemetriestudie an Europäischen Braunbären in den rumänischen Karpaten. Er liegt
als Datei \texttt{1\_bears\_RO.csv} vor und umfasst 59.997 Ortungen von 13~Tieren im
Zeitraum vom 9.~August 2011 bis zum 31.~Juli 2015.

Jede Zeile enthält neben Kennung und Zeitstempel die Koordinaten in EPSG:3844 (siehe
Abschnitt~\ref{subsec:crs}) sowie vier beschreibende Merkmale: Geschlecht, Altersklasse und
zwei Angaben zur Jahresphase. Tabelle~\ref{tab:karpaten-merkmale} zeigt deren Verteilung.

\begin{table}[htbp]
\centering
\caption{Verteilung der Merkmale im Karpaten-Datensatz ($n = 59.997$)}
\label{tab:karpaten-merkmale}
\small
\begin{tabular}{l r l r}
\toprule
\textbf{Jahresphase} & \textbf{Ortungen} & \textbf{Merkmal} & \textbf{Ortungen} \\
\midrule
Winter sleep              & 15.441 & männlich  & 36.315 \\
Den exit and reproduction & 13.199 & weiblich  & 23.682 \\
Forest fruits             & 15.091 & adult     & 41.052 \\
Hyperphagia               & 16.266 & juvenil   & 18.945 \\
\bottomrule
\end{tabular}
\end{table}

Zwei Eigenschaften dieses Bestands sind für die weitere Arbeit bedeutsam.

Erstens sind die Jahresphasen bereits Bestandteil der Rohdaten. Sie wurden nicht im Rahmen
dieser Arbeit oder einer Vorarbeit gebildet, sondern liegen als eigene Spalte vor. Die
Bezeichnungen folgen dem Jahreszyklus des Braunbären und sind damit fachlich und nicht
kalendarisch begründet.

Zweitens sind die Ortungen sehr ungleich auf die Tiere verteilt. Der am dichtesten
beobachtete Bär trägt 12.313 Ortungen bei, der am dünnsten beobachtete 1.024 - ein
Verhältnis von etwa zwölf zu eins. Auswertungen über alle Tiere hinweg werden dadurch von
wenigen Individuen dominiert. Dieser Umstand ist bei der Bewertung der in
Kapitel~\ref{chap:implementierung} vorgenommenen Messungen zu berücksichtigen.



\subsection{Kameradaten der Region Aizu}
\label{subsec:japandaten}

Der japanische Datenbestand ist strukturell anders und deutlich kleiner. Er besteht aus
20~Gerätestandorten im Aizu-Becken (Präfektur Fukushima) sowie den dort erfassten
Kameraaufnahmen. Die Standorte verteilen sich auf vier räumlich getrennte Gruppen,
erkennbar an den Kennungen der Form \texttt{001xxx} bis \texttt{004xxx}.

Insgesamt liegen 53~Aufnahmen vor, davon zwei aus Testläufen, mit denen die Bildübernahme
aus dem Postfach geprüft wurde; sie werden bei allen Auswertungen ausgeschlossen. Nach Abzug einer weiteren Aufnahme ohne Zeitstempel
verbleiben 50~verwertbare Zeitpunkte auf 18~Gerätekennungen. Nur zehn davon stehen in der
Standortliste; die übrigen acht treten erst in Aufnahmen vom November und Dezember~2025 auf
und haben keine Koordinaten. Diese Bilder sind zeitlich, nicht aber räumlich auswertbar.

Festzuhalten ist ferner, was über die Geräte \emph{nicht} gespeichert ist. Die Gerätetabelle
führt je Eintrag nur eine Kennung, den Namen des Bildverzeichnisses und die Koordinaten.
Die Bauart des Geräts oder die Art der Abschreckung ist nicht hinterlegt. Auswertungen, die
zwischen Gerätetypen unterscheiden, sind auf dieser Datengrundlage folglich nicht möglich.

\subsection{Datenhaltung im Ausgangszustand}
\label{subsec:datenhaltung}

Welches Teilsystem auf welche Ablage zugriff, zeigt Abschnitt~\ref{subsec:systemuebersicht}.
Hier geht es um die zentrale MariaDB-Instanz selbst.

Das Schema umfasste elf Tabellen: die Gerätestandorte, die von Nutzenden gezeichneten Marker
und Flächen, die Bildmetadaten, die Karpaten-Ortungen sowie einige abgeleitete Tabellen zu
Aktionsräumen und Tagesdistanzen. Drei Eigenschaften erschweren die Arbeit damit.

Erstens sind alle Zeitstempel als Zeichenketten gespeichert, nicht als Datumstyp. Ein Filter
über einen Zeitraum muss deshalb auf Text arbeiten oder erst nach dem Laden in der Anwendung
greifen. Die Datenbank kann die Ergebnismenge also nicht wirksam verkleinern.

Zweitens liegen die Koordinaten in einfachen Gleitkommaspalten. Räumliche Datentypen kommen
nicht zum Einsatz. Wer nur den sichtbaren Kartenausschnitt laden will, muss die Spalten
einzeln vergleichen.

Drittens liegen die Kameraaufnahmen in zwei Tabellen. \texttt{image\_metadata} speichert je
Aufnahme einen Dateipfad, das Bild selbst liegt im Dateisystem. \texttt{wildlife\_images}
speichert die Bilddaten direkt in der Datenbank. Der Grund dafür ist dokumentiert: Für die
Anbindung der Bildextraktion wurde die Ablage von Dateipfaden auf Bilddaten umgestellt
\cite[S.~21]{LelchatBA}, die ältere Tabelle blieb daneben bestehen. Eine Abfrage über alle
Bilder war damit in einem Schritt nicht möglich: Je nachdem, welches Teilsystem eine
Aufnahme eingetragen hatte, stand sie in der einen oder in der anderen Tabelle.

Hinzu kommt ein Fehler in den Karpaten-Ortungen. Die Tabelle \texttt{bears\_tracking}
enthielt 116.014 Datensätze, die Quelldatei aber nur 59.997 Ortungen. Der Grund liegt im
Import: Er liest zwei CSV-Dateien ein und hängt sie aneinander, obwohl beide denselben
Zeitraum abdecken und sich nur in zusätzlichen Saisonangaben unterscheiden. Die Tabelle hat
keinen eindeutigen Schlüssel aus Tier und Zeitpunkt, deshalb wurden die doppelten Ortungen
beim Einfügen nicht abgewiesen. Auf der Karte würde dieselbe Position mehrfach gezeichnet,
und abgeleitete Werte wie Schrittlängen wären falsch.

% =====================================================================
\section{Übernommene Modelle und ihre Grenzen}
\label{sec:uebernommene-modelle}

\subsection{Bewegungssimulation}
\label{subsec:simulationsmodell}

\paragraph{Aufbau des Modells.}

Das Bewegungsmodell aus \cite{BreuBA} setzt das Prinzip der Step-Selection-Funktionen um
(Abschnitt~\ref{subsec:bewegungsoekologie}). Die Kandidatenerzeugung liegt in der Klasse
\texttt{StepSelectionModel}. Das Zerlegen der realen Trajektorien in Segmente und deren
Trennung nach Jahresphasen übernimmt die Klasse \texttt{StepTrajectorySet}, die
\texttt{StepSelectionModel} dafür heranzieht.

Die Schrittlänge kommt aus einer Tabelle mit Geschwindigkeitswerten, die in der
mitgelieferten SQLite-Datei liegt. Die Tabelle hat 384~Einträge, gegliedert nach Geschlecht,
Altersklasse, Jahresphase und Stunde des Tages: zwei Geschlechter, zwei Altersklassen, vier
Jahresphasen und 24~Stunden. Je Eintrag stehen ein Mindest-, ein Höchst- und ein Medianwert
der Geschwindigkeit, wobei diese drei Werte bei 284 der 384~Einträge übereinstimmen.

Die Alternativschritte entstehen so: Für jeden realen Schritt sucht das Modell die mediane
Geschwindigkeit zu Geschlecht, Alter, Jahresphase und Stunde und rechnet daraus die
Schrittlänge. Dann bildet es $k$~Alternativen. Deren Richtung wird zufällig und
gleichverteilt über den ganzen Kreis gezogen. Je Alternative kommt ein Wahrnehmungspunkt in
200~Metern Entfernung dazu, an dem sich später Geländeeigenschaften auswerten lassen. Der
reale Schritt wird derselben Liste angefügt und als solcher markiert.

\paragraph{Ausgabe des Modells.}

Die Ausgabe ist eine Liste von Listen: je realem Schritt eine Menge von Kandidaten, jeder
mit Start- und Endpunkt, Wahrnehmungspunkt, Richtung, Länge und der Markierung, ob er der
reale Schritt ist.

Das Modell erzeugt damit die Kandidaten. Die Auswahl zwischen ihnen gehört nicht zu seinem
Umfang: Es gibt keine Bewertungsfunktion und keinen Schritt, der aus mehreren Auswahlen
einen zusammenhängenden Pfad bildet. Die als Einstiegspunkt vorgesehene Datei
\texttt{Simulation.py} ist leer.\footnote{Dateigröße 0~Byte im übernommenen Stand des
Projekts.} Vom Prinzip aus Abschnitt~\ref{subsec:bewegungsoekologie} liegt also die erste
Hälfte vor - die Menge der verfügbaren Schritte -- nicht aber die zweite.

\paragraph{Was für diese Arbeit ergänzt werden musste.}
\label{subsubsec:limitationen-simulation}

Neben der Auswahl waren vier Punkte offen.

\textbf{Richtungstreue.} Die Richtungen sind gleichverteilt über den ganzen Kreis. Die
Richtungspersistenz realer Tierbewegung (Abschnitt~\ref{subsec:bewegungsoekologie}) ist
darin nicht enthalten. Ein Pfad direkt aus diesen Kandidaten wäre ein einfacher und kein
korrelierter Zufallspfad.

\textbf{Schrittlänge.} Alle Alternativen eines Schritts haben dieselbe Länge, abgeleitet aus
dem Median. Mindest- und Höchstwert gehen nicht ein, die Streuung der Schrittlängen damit
auch nicht -- bei den oben genannten 284~Einträgen gibt es ohnehin keine.

\textbf{Zeitbezug der Schrittlänge.} Die Schrittlänge wird als Strecke einer Stunde
berechnet, unabhängig davon, wie weit die beiden zugehörigen Ortungen zeitlich
auseinanderliegen. Da rund ein Fünftel der Messintervalle länger als eine Stunde ist, fallen
die Kandidaten dort kürzer aus, als es zum Zeitabstand passt.

\textbf{Abhängigkeiten.} Eine eigene Klasse fragt Geländemerkmale über OpenStreetMap ab, ist
aber nicht an die Schrittauswahl angebunden. Außerdem lädt schon das Erzeugen eines
\texttt{StepSelectionModel}-Objekts eine Kette weiterer Klassen, die über eine
Cluster-Analyse bis zu einer kompilierten Fremdbibliothek reicht -- auch dann, wenn nur die
Schritterzeugung gebraucht wird. Lässt sich diese Bibliothek auf dem Zielsystem nicht laden,
schlägt bereits der Import fehl.

Aus diesen Punkten folgt die Entscheidung, die benötigte Logik direkt neu zu implementieren
und nicht das vorhandene Objekt einzubinden (Abschnitt~\ref{sec:uebertragung}).

\subsection{Bildextraktion}
\label{subsec:bildextraktion}

Die dritte Vorarbeit \cite{LelchatBA} übernimmt neue Kameraaufnahmen automatisch. Die
Kameras schicken ihre Aufnahmen als E-Mail-Anhang. Das System holt diese Anhänge ab, wertet
sie aus und trägt sie in die Datenbank ein.

\paragraph{Ablauf.}

Das System läuft als eigener Prozess, getrennt von der Kartenanwendung. Es fragt ein
Postfach über die Microsoft Graph API ab und meldet sich dabei per OAuth~2.0 als Anwendung
an, nicht als Benutzer. Die Abfrage läuft in einer Dauerschleife alle zehn Sekunden. Je
Durchlauf holt das System bis zu 50~ungelesene Nachrichten. Der Filter auf einen festen
Absender greift danach im Programm, nicht schon in der Abfrage.

Aus jedem Anhang im JPEG-Format gewinnt das System Information auf zwei Wegen:

\begin{itemize}
    \item \textbf{Aus dem Bild} liest die Texterkennung \texttt{easyOCR} den eingeblendeten
    Text. Daraus holt ein regulärer Ausdruck den Zeitstempel.
    \item \textbf{Aus dem Dateinamen} kommt die Gerätekennung. Das System kennt dafür zwei
    Namensformate. Trägt eine Aufnahme keinen eingeblendeten Zeitstempel -- was vorkommt
    \cite[S.~25]{LelchatBA} --, stammt die Zeit ersatzweise aus dem Dateinamen. Das gelingt
    nur bei einem der beiden Formate, weil nur dieses ein Datum in auswertbarer
    Schreibweise enthält.
\end{itemize}

Danach schreibt das System die Aufnahme mit Gerätekennung und Zeitstempel in die Datenbank.
Ein Merkmal hält fest, ob beides bestimmt werden konnte. Ist die Gerätekennung in der
Gerätetabelle noch nicht vorhanden, legt das System dort einen neuen Eintrag an. Dieser
Eintrag erhält zunächst die Koordinaten 0/0.

\paragraph{Einordnung der Artbestimmung.}

Die Tierart wird nicht von diesem System bestimmt. An den Kameras im Gelände stehen
Raspberry-Pi-Rechner, die die Aufnahmen auswerten, Rahmen mit Klassenname und
Konfidenzwert in das Bild zeichnen und erst dann versenden \cite[S.~8]{LelchatBA}. Die
Kamera selbst blendet nur Gerätenummer, Zeitpunkt und Temperatur ein. Die Texterkennung
liest also die Beschriftung dieser Rahmen: Steht darin das Wort \enquote{bear}, wird die
Aufnahme übernommen, andernfalls verworfen \cite[S.~20]{LelchatBA}.

Die Artangabe stammt damit aus einem vorgeschalteten Modell, dessen Fehlerrate nicht bekannt
ist; Fehlerkennungen sind dokumentiert \cite[S.~24]{LelchatBA}. Sie wird, wie in
Abschnitt~\ref{sec:zielsetzung} festgelegt, als gegeben vorausgesetzt und in dieser Arbeit
weder geprüft noch verbessert.

\paragraph{Zwei Punkte für die Weiterarbeit.}

\begin{itemize}
    \item \textbf{Schreibrechte auf dem Postfach.} Nach der Verarbeitung versucht das
    System, die Nachricht als gelesen zu markieren, denn die Abfrage holt nur ungelesene
    Nachrichten. Das setzt eine Schreibberechtigung auf dem Postfach voraus. Diese war zur
    Zeit der Vorarbeit nicht erteilt; der Quelltext hält das fest.\footnote{Kommentar in
    \texttt{WildlifePictureExtractor.py} im Anschluss an die Markierungsanfrage.} Ohne die
    Berechtigung bleibt eine Nachricht ungelesen und wird im nächsten Durchlauf erneut
    verarbeitet. Der Verarbeitungsstand liegt damit im Postfach und nicht in der Anwendung.

    \item \textbf{Keine Rückmeldung an die Kartenanwendung.} Extraktion und Kartenanwendung
    sind nur über die gemeinsame Datenbank verbunden. Die Kartenanwendung erfährt nicht,
    dass neue Aufnahmen eingetroffen sind; Nutzende müssten selbst nachsehen.
\end{itemize}

% =====================================================================
\section{Anforderungen}
\label{sec:anforderungen}

Aus den vorangegangenen Abschnitten ergeben sich die folgenden Anforderungen. Sie sind
nummeriert, damit in Kapitel~\ref{chap:implementierung} und
Kapitel~\ref{chap:diskussion} auf sie verwiesen werden kann.

\subsection{Funktionale Anforderungen}

\begin{description}
    \item[F1 - Gemeinsame Karte.] Alle Datenbestände sind auf einer einzigen Karte
    darstellbar. Einzelne Datenebenen lassen sich unabhängig voneinander ein- und
    ausblenden.

    \item[F2 - Filterung.] Die dargestellten Daten lassen sich nach Tier, Zeitraum und
    Jahresphase filtern.

    \item[F3 - Eigene Marker und Flächen.] Nutzende können auf der Karte Marker und
    Flächen anlegen, dauerhaft speichern und wieder entfernen - und zwar unabhängig
    davon, welche Datenebene betrachtet wird.

    \item[F4 - Zeitlicher Verlauf.] Der zeitliche Ablauf einer Bewegung ist abspielbar,
    statt nur als abgeschlossene Linie dargestellt zu werden.

    \item[F5 - Gegenüberstellung von Modell und Beobachtung.] Zu einer realen Trajektorie
    lassen sich die vom Modell erzeugten Alternativen für dasselbe Tier, denselben Tag und
    dieselben Zeitpunkte anzeigen.

    \item[F6 - Erzeugung zusammenhängender Pfade.] Das Bewegungsmodell erzeugt
    fortlaufende Pfade. Dazu gehören eine Auswahl zwischen den Kandidatenschritten, eine
    an realen Daten geschätzte Richtungstreue und eine Bewertung des Geländes.

    \item[F7 - Statistischer Vergleich.] Die Übereinstimmung zwischen simulierter und
    beobachteter Tagesaktivität wird als Kennzahl mit Unsicherheitsangabe berechnet.

    \item[F8 - Einsicht in die Rohdaten.] Zu jeder Kartendarstellung sind die
    zugrunde liegenden Datensätze einsehbar.

    \item[F9 - Sichtbarkeit neuer Aufnahmen.] Seit der letzten Betrachtung neu
    eingetroffene Kameraaufnahmen werden in der Anwendung kenntlich gemacht.
\end{description}

\subsection{Nicht-funktionale Anforderungen}

\begin{description}
    \item[N1 - Darstellungsleistung.] Die Darstellung bleibt auch dann nutzbar, wenn
    mehrere Datenebenen mit zusammen mehreren zehntausend Objekten gleichzeitig sichtbar
    sind.

    \item[N2 - Keine schwergewichtigen Abhängigkeiten.] Die Anwendung startet und läuft
    auch dann, wenn eine kompilierte Fremdbibliothek der Vorarbeit nicht ladbar ist.

    \item[N3 - Einheitliche Datenhaltung.] Die Daten werden über eine gemeinsame
    Zugriffsschicht gelesen. Lokale Dateien dienen allenfalls als Rückfallebene, nicht als
    eigenständige Quelle.

    \item[N4 - Nachvollziehbare Fehler.] Ist eine Datenquelle nicht erreichbar, wird das
    kenntlich gemacht. Eine leere Ergebnismenge und eine fehlgeschlagene Verbindung sind
    für Nutzende unterscheidbar.

    \item[N5 - Reproduzierbarkeit.] Simulationsläufe und daraus abgeleitete Kennzahlen
    sind bei gleicher Parametrisierung wiederholbar.

    \item[N6 - Übertragbare Konfiguration.] Zugangsdaten und umgebungsabhängige
    Einstellungen sind nicht fest im Quelltext hinterlegt.

    \item[N7 - Einheitlicher Bildbestand.] Alle Kameraaufnahmen sind über eine einzige
    Abfrage erreichbar, unabhängig davon, welches Teilsystem sie eingetragen hat.

    \item[N8 - Eindeutigkeit der Ortungen.] Jede Ortung eines Tieres zu einem bestimmten
    Zeitpunkt wird genau einmal dargestellt und genau einmal in abgeleitete Kennzahlen
    einbezogen, auch wenn der Datenbestand mehrfach vorliegende Datensätze enthält.
\end{description}

% ====================================================================
\chapter{Implementierung}
\label{chap:implementierung}

\section{Kartenkomponente und Interaktionskonzept}
\label{sec:kartenvisualisierung}

\subsection{Architekturentscheidung: statische gegen bidirektionale Komponente}
\label{subsec:architekturentscheidung}

Streamlit bietet zwei Wege, eigenes HTML und JavaScript einzubinden
(Abschnitt~\ref{par:komponenten}). Beide kamen in Frage, und die Wahl bestimmt den Aufbau des
ganzen Kapitels.

\paragraph{Die bidirektionale Komponente.} \texttt{declare\_component()} legt einen Rückkanal
an. Der JavaScript-Teil ruft \texttt{Streamlit.setComponentValue()} auf, und Python erhält
diesen Wert als Rückgabewert. Eine gezeichnete Fläche oder ein angeklickter Punkt wäre damit
sofort als Python-Objekt verfügbar.

Dieser Weg verlangt aber ein eigenes Frontend-Projekt: ein Verzeichnis mit
\texttt{package.json}, Paketen aus dem npm-Ökosystem und einem Übersetzungsschritt, dessen
Ergebnis vor dem Start der Anwendung vorliegen muss. Auf dem Zielsystem bräuchte man neben
Python auch Node.js, und jede Änderung an der Karte würde einen neuen Übersetzungslauf nötig
machen.

\paragraph{Die statische Komponente.} \texttt{components.html()} übergibt eine Zeichenkette
mit HTML an einen \texttt{iframe} und liefert nichts zurück. Es gibt keinen
Übersetzungsschritt, keine zweite Laufzeitumgebung und keine Paketverwaltung. Die Komponente
liegt in einer einzigen Python-Datei, und MapLibre~GL~JS wird zur Laufzeit über ein Content
Delivery Network geladen.\footnote{Eingebunden ist Version 3.6.2 über \texttt{unpkg.com}, in
allen Kartenkomponenten dieselbe.}

\paragraph{Entscheidung.} Gewählt wurde \texttt{components.html()}. Der Grund: Das System soll
an der Hochschule weitergeführt werden. Eine Anwendung, die man mit \texttt{streamlit run}
startet und deren Kartenlogik in einer lesbaren Python-Datei steht, ist dafür leichter
zugänglich als eine Anwendung mit getrenntem Frontend-Build. Dazu kommt: Der Rückkanal nützt
hier wenig. Fast alles, was man auf der Karte tut, betrifft nur die Darstellung. Ebenen
schalten, Stil wechseln, Kamera bewegen -- das muss Python gar nicht erreichen.

Der Preis ist der fehlende Rückkanal für alles, was die Daten verändert, also für neue Marker
und Gebiete. Wie er bezahlt wurde, zeigt Abschnitt~\ref{sec:interaktion}.

\subsection{Aufbau der Kartenkomponente}
\label{subsec:aufbau-karte}

Die Komponente ist eine Funktion namens \texttt{render\_maplibre()}. Sie bekommt die Daten,
baut daraus ein HTML-Dokument und gibt es über \texttt{components.html()} an Streamlit
weiter.

\paragraph{Datenübergabe.} Die Daten kommen als GeoJSON an, in fünf Gruppen: die japanischen
Gerätestandorte mit den dort angelegten Markern und Gebieten, die Karpaten-Trajektorien, die
simulierten Bewegungen und zwei externe Vergleichsbestände
(Abschnitt~\ref{sec:externe-quellen}). Sie werden direkt in das HTML-Dokument geschrieben.
Die Karte fragt deshalb selbst nichts bei der Datenbank ab. Fehlt eine Gruppe, bleibt ihre
Ebene leer und der Rest der Karte funktioniert weiter. Abbildung~\ref{fig:hauptkarte} zeigt die Karte mit mehreren eingeschalteten Ebenen.

\begin{figure}[htbp]
\centering
\includegraphics[width=0.4\linewidth]{karte-karpaten.png}
\caption[Hauptkarte mit mehreren Datenebenen]{Hauptkarte mit eingeschalteter
  Geländedarstellung. Sichtbar sind die GPS-Pfade der Karpatenbären und die simulierten
  Bewegungen. Links die Ebenenliste mit der Anzahl der Objekte je Gruppe. Im Ausgangszustand
  lagen diese Bestände auf zwei getrennten Seiten
  (Abbildung~\ref{fig:ausgangszustand-karten}).}
\label{fig:hauptkarte}
\end{figure}

\paragraph{Quellen und Ebenen.} Jede Gruppe bekommt eine eigene Datenquelle. Darauf liegen
die sichtbaren Ebenen: Punkte als \texttt{circle}, Pfade als \texttt{line}, Gebiete als
\texttt{fill} mit Umrandung (Abschnitt~\ref{subsec:maplibre}). Ändern sich die Daten, wird
nur die Quelle ausgetauscht und die Ebenen bleiben. Die Animation nutzt das und tauscht je
Einzelbild nur die Daten aus.

Neue Ebenen liegen über allen bestehenden und würden die Ortsnamen der Grundkarte verdecken.
Sie werden deshalb darunter eingefügt.

\paragraph{Stilwechsel.} Wechselt man die Grundkarte, lädt MapLibre den Stil komplett neu.
Die eigenen Datenebenen gehen dabei verloren, weil sie Teil des Stils sind. Die Komponente
wartet deshalb, bis der neue Stil geladen ist, und legt die Ebenen danach wieder an.

\paragraph{Grundkarten und Gelände.} Zur Auswahl stehen sieben Grundkarten, vier Vektor- und
drei Rasterstile. Voreingestellt ist das Luftbild. Nur darauf sieht man, ob ein Pfad über
Wald, Lichtung oder Siedlung führt. Für jede Rasterquelle ist hinterlegt, bis zu welcher
Zoomstufe der Dienst Kacheln liefert; ohne diese Angabe fordert die Karte Kacheln an, die es
nicht gibt.

Dazu kann man eine dreidimensionale Ansicht einschalten. Sie beruht auf einem Höhenmodell in
Kachelform (Abschnitt~\ref{subsec:hoehenmodelle}) und erzeugt zwei Dinge: eine Schattierung,
die die Geländeform zeigt, und eine Anhebung der Karte selbst. Ist die Höhenquelle nicht
erreichbar, schaltet sich der Schalter von selbst aus und zeigt einen Hinweis. Die Karte
läuft dann zweidimensional weiter. Dasselbe Verhalten gibt es an anderen Stellen im System
(Abschnitt~\ref{sec:datenhaltung}).


 

\subsection{Interaktion auf der Karte}
\label{sec:interaktion}

Weil die Komponente statisch ist, fließt vom Browser kein Wert nach Python zurück.

\paragraph{Was im Browser bleibt.} Die meisten Aktionen ändern nur die Anzeige: Ebenen ein-
und ausblenden, den Stil wechseln, die Karte bewegen, ein Objekt anklicken. Das passiert
alles im Browser, ohne Rerun (Abschnitt~\ref{subsec:streamlit}).

Beim Einblenden springt die Karte zur Ebene. Punkte bei (0,0) lässt sie dabei weg, denn
Geräte ohne Standort stehen mit diesen Werten in der Datenbank
(Abschnitt~\ref{subsec:japandaten}).

\paragraph{Der Kartenzustand.} Wendet man einen Filter an, baut Streamlit die Seite neu auf,
und die Karte würde auf ihre Ausgangsansicht zurückspringen. Die Komponente schreibt den
Kartenzustand deshalb in die Adresszeile: Mittelpunkt, Zoomstufe, Stil, Zustand des
3D-Schalters und die sichtbaren Ebenen. Das ändert nur die Adresse und lädt die Seite nicht
neu. Weil eine Kartenbewegung viele Ereignisse auslöst, wird höchstens alle 400~Millisekunden
geschrieben. Beim nächsten Aufbau liest Python diese Werte aus der Adresse und gibt sie an
die Komponente zurück, die an derselben Stelle startet.

Verbieten die Sicherheitsregeln des Browsers den Zugriff auf das umgebende Fenster, bleibt
die Karte nutzbar und verliert nur diese Zustandsübernahme.

\paragraph{Neue Marker und Gebiete.} Nur diese Aktion geht an Python, und sie nutzt denselben
Weg (Quelltext~\ref{lst:rueckweg}). Die Karte hat dafür ein eigenes Bedienfeld mit zwei Modi.
Im Markermodus setzt ein Klick einen Punkt, im Gebietsmodus umrandet man eine Fläche durch
mehrere Klicks. Die Eingabe erscheint währenddessen als eigene Vorschauebene. Einzelne Punkte
lassen sich zurücknehmen, und bis zum Speichern kann man alles abbrechen. Der Speichern-Knopf
bleibt gesperrt, bis die Eingabe gültig ist: bei einem Gebiet mindestens drei Punkte und ein
Name.

\begin{lstlisting}[language=JavaScript, caption={Derselbe Weg über die Adresszeile, einmal
ohne und einmal mit Neuladen. Gekürzt aus \texttt{saveStateToUrl()} und
\texttt{saveDrawing()}.}, label={lst:rueckweg}]
// Kamera und Ebenen: nur die Adresse aendern, kein Neuladen
parentWin.history.replaceState(null, '', newUrl);

// Neues Gebiet: Werte anhaengen, dann echte Navigation
url.searchParams.set('new_polygon_name', nameInput.value.trim());
url.searchParams.set('new_polygon_coords', JSON.stringify(drawState.polygonPoints));
parentWin.location.href = url.toString();
\end{lstlisting}

Beim Speichern schreibt die Komponente die Geometrie in die Adresszeile und lädt die Seite
wirklich neu. Python findet die Werte im nächsten Durchlauf, legt den Eintrag in der
Datenbank an und löscht die Werte sofort aus der Adresse. Sonst würde beim nächsten Laden
ein zweiter Eintrag entstehen. Danach wird der Zwischenspeicher der japanischen Daten
geleert, weil er sonst den alten Stand behält und der neue Marker nicht erscheint
(Abschnitt~\ref{subsec:streamlit}).


\paragraph{Löschen, Filter und Rohdaten.} Gelöscht wird unterhalb der Karte in einer
Tabelle, also mit normalen Streamlit-Elementen, deren Wert Python ohnehin bekommt. Der Knopf
bleibt gesperrt, bis man zusätzlich bestätigt. Die beiden teuren Filter für Karpaten und
Simulation stecken jeweils in einem Formular: Die Eingaben werden gesammelt und erst mit
\enquote{Anwenden} gemeinsam wirksam. Sonst würde jede einzelne Änderung einen Neuaufbau
samt Datenbankabfrage auslösen. Unter der Karte stehen dieselben Daten zusätzlich als
Tabelle, nach Quelle in Reiter getrennt.



\section{Integration der Teilsysteme und Vereinheitlichung der Datenhaltung}
\label{sec:integration}
\subsection{Anbindung der Bewegungssimulation}
  \label{subsec:anbindung-simulation}

      Die Bewegungssimulation sollte als eigene Ebene in das System eingebunden
  werden, sodass sich reale und simulierte Bewegungen in derselben Karte
  vergleichen lassen. Das übernommene System liest dafür GPS-Daten aus einer
  SQLite-Datenbank und bildet daraus Trajektoriensegmente: Ein Segment endet bei
  einer Zeitlücke von mehr als 1,5~Stunden und wird zusätzlich an Saisongrenzen
  getrennt. Zu jedem realen Schritt erzeugt das Modell dann einen Fächer aus
  Alternativen, deren Länge aus einer Tabelle mit Mediangeschwindigkeiten folgt
  und deren Richtung gleichverteilt über den Vollkreis gezogen wird.

  \paragraph{Vorgehen bei der Anbindung.} Für die Anbindung wurde eine eigene
  Integrationsschicht geschrieben. Sie stellt nach außen nur zwei Funktionen
  bereit: eine liefert die auswählbaren Bären mit ihren Saisons, die andere
  liefert zu einer Auswahl fertiges GeoJSON. Damit hat die Simulation dieselbe
  Schnittstelle wie alle anderen Datenbestände, und die Karte behandelt sie wie
  eine gewöhnliche Ebene (Abschnitt~\ref{subsec:aufbau-karte}). Das Modell selbst
  bleibt unverändert.

  \paragraph{Nur die benötigten Teile ausführen.} Der Einstiegspunkt des
  Originals führt drei Schritte nacheinander aus: Segmente bilden, nach Saison
  trennen und anschließend geometrische Mediane der Cluster bestimmen. Der dritte
  Schritt zieht eine Clusteranalyse nach sich und bringt zwei zusätzliche Pakete mit, \texttt{hdbscan} und \texttt{scipy}, ohne die sich das Modul nicht laden lässt.

Die Schritterzeugung braucht diesen Schritt aber nicht: Das dort gesetzte Attribut   \texttt{local\_geometric\_median} wird an keiner Stelle der Schritterzeugung
  gelesen. Die Integrationsschicht bildet die Segmente deshalb direkt aus der
  Datenbank, nach denselben vier Regeln wie das Original, und lässt die
  Clusteranalyse aus. Die Schritterzeugung braucht nur die
  Geschwindigkeitstabelle und die geometrischen Hilfsfunktionen.

  \paragraph{Zustand über mehrere Aufrufe.} Das Original hält eine offene
  Datenbankverbindung. Streamlit führt das Skript bei jeder Interaktion erneut
  aus, und zwar aus wechselnden Threads, während eine SQLite-Verbindung an
  genau einen Thread gebunden ist. Die Geschwindigkeitstabelle wird deshalb
  einmal vollständig gelesen und danach aus dem Arbeitsspeicher bedient.

  \paragraph{Schrittlänge und Messintervall.} Die Länge eines Schrittes ergibt
  sich im Original aus der Mediangeschwindigkeit multipliziert mit einer Stunde.
  Die Karpaten-Ortungen liegen aber nicht durchgehend eine Stunde auseinander:
  Bei rund 80~\% der Zeitabstände ist es etwa eine Stunde, bei etwa 13~\% etwa zwei Stunden, bei knapp 7~\% mehr. Betrachtet man einen einzelnen
  Schritt für sich, spielt das keine Rolle. Sobald mehrere Schritte zu einem Pfad
  zusammengesetzt werden, schon: Bei einer Zweistundenlücke wären die Alternativen nur halb so lang wie der reale Schritt, und der simulierte Pfad bliebe hinter dem realen Tier zurück. Die Schrittlänge kann deshalb mit der tatsächlich verstrichenen Zeit skaliert
werden. Auf der Hauptkarte ist das abgeschaltet, weil die Segmente dort ohnehin
an Lücken über 1,5~Stunden getrennt sind; in der Animation ist es aktiv.

  Zusätzlich erhielt die Schritterzeugung einen eigenen Zufallsgenerator. Mit einem festen Startwert liefert ein Lauf wieder dasselbe
  Ergebnis, was für die automatisierten Tests und für reproduzierbare Abbildungen Voraussetzung ist.

  \paragraph{Prüfung der Gleichwertigkeit.} Da Teile des Ablaufs nachgebaut
  wurden, musste gezeigt werden, dass dabei nichts verändert wurde. Dafür wurden
  beide Wege nebeneinander ausgeführt und ihre Ausgaben verglichen
  (Tabelle~\ref{tab:gleichwertigkeit}). Verglichen wurden einmal die gebildeten
  Segmente, Punkt für Punkt, und einmal die erzeugten Schritte bei identischem
  Zufallsstrom.

  \begin{table}[htbp]
  \centering
  \caption{Vergleich der nachgebauten Teile mit dem ursprünglichen Ablauf.}
  \label{tab:gleichwertigkeit}
  \begin{tabular}{lll}
  \toprule
  Geprüfter Teil & Umfang & Ergebnis \\
  \midrule
  Segmentbildung   & 1\,258 Segmente, 47\,143 Schritte & alle identisch \\
  Schritterzeugung & 1\,808 Schritte, je 10 Alternativen & alle identisch \\
  \bottomrule
  \end{tabular}
  \end{table}

  Beide Vergleiche stimmen vollständig überein. Die Anbindung verändert das
  Modell also nicht, sie führt es nur in einer anderen Umgebung aus.

  \paragraph{Ergebnis und offener Punkt.} Auswählbar ist eine Kombination aus Bär und Saison. Dazu zeigt die Karte den
realen Pfad und den Fächer der Alternativen je Schritt. Lässt sich das Modell
  nicht laden, zeigt die Ebene die reinen GPS-Spuren aus der Datenbank, damit die
  Karte nutzbar bleibt.

  Offen bleibt die Frage, welche der Alternativen gewählt werden soll. Solange die
Richtung gleichverteilt gezogen wird, entsteht kein Pfad, der mit dem Verhalten realer Bären
  vergleichbar wäre.

\subsection{Anbindung der Bildextraktion}
\label{subsec:anbindung}

Für die Anbindung der Bildextraktion standen zwei Architekturansätze zur Auswahl. Der erste Ansatz verlagert die Auswertung direkt in die Anwendung und fragt das Postfach bei jedem Seitenaufruf ab. Der zweite Ansatz führt die Extraktion als eigenständigen Hintergrundprozess aus und verbindet beide Komponenten über eine gemeinsame Datenbank.

Es wurde der zweite Weg gewählt. Dafür sprechen drei wesentliche Gründe:
Erstens treffen neue Aufnahmen kontinuierlich ein, unabhängig davon, ob die Benutzeroberfläche geöffnet ist. Eine Auswertung nur beim Seitenaufruf würde Daten erst verspätet erfassen. Zweitens führt Streamlit das Skript bei jeder Benutzerinteraktion neu aus; jeder Klick würde somit einen erneuten Postfachzugriff auslösen. Drittens kann die Auswertung bei vielen Nachrichten mehr Zeit in Anspruch nehmen als der Seitenaufbau, was die Oberfläche blockieren würde. Die getrennte Architektur hält beide Teile im passenden Takt: Die Extraktion arbeitet fortlaufend im Hintergrund, während die Oberfläche lediglich den aktuellen Stand aus der Datenbank ausliest.

Die Extraktion wird in der Seitenleiste der Oberfläche über eine Statusanzeige sichtbar, die neu eingetroffene Aufnahmen meldet. Hierbei wurden zwei Entscheidungen getroffen. Die erste betrifft die Abfragehäufigkeit: Um die Datenbank bei wiederholten Streamlit-Interaktionen nicht unnötig zu belasten, wird die Abfrage für 30~Sekunden zwischengespeichert (Caching). Dieser Wert stellt einen optimalen Kompromiss zwischen hoher Aktualität und geringer Systemlast dar. Die zweite Entscheidung betrifft die Speicherung des Prüfzeitpunkts. Würde dieser nur im flüchtigen Sitzungszustand halten, gälten nach jedem Neustart alle Bilder erneut als neu. Der Zeitpunkt wird daher in einer Statustabelle der Datenbank abgelegt und bleibt somit über Neustarts und Sitzungswechsel hinweg erhalten.

\paragraph{Einschränkung beim Postfachzugriff.} 
Für den dauerhaften Betrieb sollten verarbeitete Nachrichten im Postfach als gelesen markiert werden, um Mehrfachverarbeitungen zu vermeiden. Die Anwendung bietet hierfür zwei Anmeldewege. Die benutzerbasierte Anmeldung scheitert an einer Richtlinie für bedingten Zugriff der Hochschule, die für dieses Verfahren keine Token ausstellt. Für einen unbeaufsichtigten Hintergrunddienst ist dieser Weg ohnehin ungeeignet, da er eine personelle Interaktion erfordert. Vorgesehen ist daher der zweite Weg über eine anwendungsbasierte Anmeldung. Die erforderliche Berechtigung wurde beantragt und vom Informationssicherheitsbeauftragten genehmigt. Da die Tokenausstellung durch das Rechenzentrum bei Abgabe noch ausstand, weist die Microsoft-Graph-Schnittstelle Anfragen aktuell mit dem Statuscode 403 ab. Die Zuordnung bereits ausgewerteter Nachrichten erfolgt deshalb vorübergehend anwendungsseitig. Sobald die Berechtigung vorliegt, kann die Funktion ohne Codeänderungen genutzt werden.

\subsection{Vereinheitlichung der Datenhaltung}
\label{sec:datenhaltung}

Vor der Zusammenführung lagen die Kamerabilder an zwei Stellen. Die
Kartenanwendung hielt in der Tabelle \texttt{image\_metadata} nur Verweise auf
Dateien im Dateisystem; die Bildextraktion legte ihre Aufnahmen in einer eigenen
Tabelle als Bilddaten ab. Beide Bestände beschrieben denselben Sachverhalt,
waren aber unterschiedlich aufgebaut, sodass sich der Bildbestand nicht in einer
Abfrage auswerten ließ. Hinzu kam, dass Verweise auf Dateien nur auf dem Rechner
gültig sind, auf dem der Ordner liegt. Für ein System, das an der Hochschule
weiterbetrieben werden soll, ist das eine Einschränkung.

\paragraph{Entwurf der gemeinsamen Tabelle.} Die neue Tabelle speichert die
Bilddaten selbst, nicht den Pfad. Damit ist ein Eintrag unabhängig vom
Dateisystem des jeweiligen Rechners, und beide Zugänge lassen sich gleich
behandeln: Aufnahmen aus dem E-Mail-Abruf liegen ohnehin nur als Daten vor und
hätten sonst eigens als Datei abgelegt werden müssen. Der Preis dafür ist eine
deutlich größere Datenbank, was für den hier vorliegenden Umfang vertretbar ist.

Woher ein Eintrag stammt, hält eine eigene Spalte fest, mit den Werten für
E-Mail-Abruf und Dateisystem. Die Herkunft bleibt damit erhalten, ohne dass der
Bestand dafür auf zwei Tabellen aufgeteilt werden muss. Der Zeitstempel wird als
Datums- und Zeitwert geführt und nicht mehr als Zeichenkette, sodass sich
Zeiträume in der Abfrage selbst eingrenzen lassen. Indizes bestehen auf den
Spalten, nach denen die Oberfläche filtert: Gerät, Bildtyp, Herkunft und
Zeitpunkt.

\paragraph{Überführung der Bestände.} Die Überführung lief in einem eigenen
Schritt und las beide Quellen nacheinander ein: die Ordner im Dateisystem und
die Tabelle der Bildextraktion. Ein Eintrag ist über Gerät, Dateiname und
Herkunft eindeutig; diese Kombination ist in der Tabelle als eindeutiger
Schlüssel hinterlegt. Zusammen mit einem Einfügen, das vorhandene Einträge
überspringt, lässt sich die Überführung beliebig oft wiederholen, ohne dass
Dubletten entstehen. Das war nötig, weil die Bildextraktion während der
Bearbeitung weiterlief und laufend neue Aufnahmen hinzukamen.

Die beiden alten Tabellen wurden nicht gelöscht. So blieb der Ausgangszustand
zum Vergleich erhalten, und ein Rückschritt wäre jederzeit möglich gewesen.

\paragraph{Ergebnis.} Alle Zugriffe der Anwendung auf Bilder richten sich gegen
diese eine Tabelle, sowohl die Anzeige der Bilder zu einem Gerät als auch die
Zählung je Gerät für die Karte. Damit ist Anforderung~N7 erfüllt: Der
Bildbestand ist über eine einzige Abfrage erreichbar, unabhängig davon, welches
Teilsystem einen Eintrag angelegt hat.

\begin{figure}[htbp]
\centering
\begin{tikzpicture}[
  font=\footnotesize,
  tab/.style={draw, inner sep=0pt},
  neu/.style={draw, line width=1pt, inner sep=0pt},
  pfeil/.style={-{Latex[length=2mm]}, line width=0.7pt},
  notiz/.style={font=\scriptsize\itshape, align=left}
]

\newcommand{\tabelle}[2]{%
  \begin{minipage}{42mm}
  \vspace{1.3mm}
  \centering\texttt{#1}\\[1mm]
  \hrule
  \vspace{1.3mm}
  \hspace*{2.5mm}\begin{minipage}{37mm}#2\end{minipage}
  \vspace{1.3mm}
  \end{minipage}}

\newcommand{\tabelleB}[2]{%
  \begin{minipage}{48mm}
  \vspace{1.3mm}
  \centering\texttt{#1}\\[1mm]
  \hrule
  \vspace{1.3mm}
  \hspace*{2.5mm}\begin{minipage}{43mm}#2\end{minipage}
  \vspace{1.3mm}
  \end{minipage}}

\node[tab] (t1) at (0,0) {\tabelle{image\_metadata}{%
  device\_id\\ filename\\ \textbf{image\_path} \textit{(Verweis)}\\ image\_type\\
  timestamp \textit{(Text)}\\ parsed\_successfully}};

\node[tab] (t2) at (0,-4.3) {\tabelle{wildlife\_images}{%
  device\_id\\ filename\\ \textbf{image\_bytes} \textit{(Bilddaten)}\\ image\_type\\
  timestamp \textit{(Text)}\\ parsed\_successfully}};

\node[neu] (t3) at (9.0,-2.15) {\tabelleB{bears\_images\_unified}{%
  device\_id\\ filename\\ \textbf{image\_bytes} \textit{(Bilddaten)}\\ image\_type\\
  \textbf{source} \textit{(Herkunft)}\\
  \textbf{timestamp} \textit{(Datumswert)}\\ parsed\_successfully}};

\draw[pfeil] (t1.east) -- (t3.west);
\draw[pfeil] (t2.east) -- (t3.west);

\node[notiz, anchor=north west] at ($(t3.south west)+(0,-0.3)$)
  {eindeutig über \texttt{(device\_id, filename, source)}\\
   Indizes auf Gerät, Bildtyp, Herkunft und Zeitpunkt};

\node[notiz, anchor=north west] at ($(t2.south west)+(0,-0.3)$)
  {bleiben als Ausgangszustand erhalten};

\end{tikzpicture}
\caption{Zusammenführung der beiden Bildtabellen. Hervorgehoben sind die
Spalten, in denen sich die Bestände unterschieden oder die neu hinzugekommen
sind.}
\label{fig:datenhaltung}
\end{figure}



\section{Bewegungsanimation}
\label{sec:animation}

\subsection{Synchronisierte Darstellung echter und simulierter Bewegung}
\label{subsec:synchron}

Eine Trajektorie als fertige Linie auf der Karte zeigt, \emph{wo} ein Tier war, aber nicht,
\emph{wann} es dort war und wie schnell es sich bewegt hat. Für die Beurteilung des
Bewegungsmodells ist gerade das entscheidend, denn das Modell erzeugt Schritte je Zeitpunkt.
Die Animationskomponente macht den zeitlichen Verlauf deshalb abspielbar.

\paragraph{Die Forderung nach Synchronität.} Entscheidend ist dabei nicht die Animation an
sich, sondern dass die simulierte Bewegung sich auf \emph{dieselbe} Grundlage bezieht wie die
reale: dasselbe Tier, denselben Tag und dieselben Zeitpunkte. Wären beide getrennt
auswählbar, ließen sich unversehens zwei inhaltlich unverbundene Datensätze nebeneinander
stellen, und der Vergleich wäre wertlos.

Die Synchronität wird deshalb nicht der Bedienung überlassen, sondern im Aufbau erzwungen:
Die Kandidatenschritte werden aus den realen Punkten desselben Tages erzeugt, und jeder
simulierte Punkt übernimmt den Zeitstempel des realen Schritts, zu dem er gehört. Beide
Darstellungen laufen damit zwangsläufig über dieselbe Zeitachse.

\paragraph{Zwei Formen der Simulation.} Die Komponente zeigt die Modellausgabe in zwei
Formen, die unterschiedliche Fragen beantworten und gleichzeitig sichtbar sein können.

\begin{wrapfigure}{r}{0.48\textwidth}
    \centering
    \includegraphics[width=0.46\textwidth]{Fächer.png}
    \caption{Der Fächer als Möglichkeitsraum.}
    \label{fig:faecher}
\end{wrapfigure}

Der \emph{Fächer} zeigt an jedem realen Schritt sämtliche Kandidaten, die das Modell dort
erzeugt hat (siehe Abbildung~\ref{fig:faecher}). Er ist kein Pfad, sondern der Möglichkeitsraum eines \emph{einzelnen} Schritts.

Die \emph{Realisierungen} sind durchgehende Pfade: Je Schritt wird ein Kandidat gezogen und
auf die aktuelle simulierte Position angewendet, nicht auf die reale. Nur so entsteht ein
fortlaufender Weg und nicht wieder ein Fächer. Dies ergibt den zweiten Bären, der neben dem
realen herläuft.

\paragraph{Warum mehrere Pfade gleichzeitig.} Standardmäßig werden mehrere Realisierungen
gezeichnet. Das ist eine bewusste Entscheidung gegen eine naheliegende, aber irreführende
Darstellung: Eine einzelne Linie neben der realen Route würde unweigerlich als
\enquote{das Modell sagt diesen Weg voraus} gelesen. 

\begin{wrapfigure}{r}{0.48\textwidth}
    \centering
    \includegraphics[width=0.46\textwidth]{mehrereSimuliertePfade.png}
    \caption{Mehrere simulierte Realisierungen zur Streuung.}
    \label{fig:mehrere_pfade}
\end{wrapfigure}

Ein Modellpfad ist jedoch eine
Zufallsziehung; ein erneuter Lauf ergibt einen anderen Weg. Mehrere auseinanderlaufende
Linien vermitteln diesen Charakter unmittelbar (vgl. Abbildung~\ref{fig:mehrere_pfade}), und aussagekräftig ist ihre Streuung, nicht
die einzelne Linie. Ein Hinweis in der Oberfläche benennt das zusätzlich ausdrücklich.

\paragraph{Gewichtete Auswahl statt Gleichverteilung.} Innerhalb eines Schritts sind die
Kandidatenrichtungen gleichverteilt, wie Abschnitt~\ref{subsec:simulationsmodell} für die
Vorarbeit dargelegt hat. Würde unter ihnen ebenfalls gleichverteilt gezogen, entstünde ein
Pfad ohne jeden Bezug zum Drehverhalten realer Tiere. Die Kandidaten werden deshalb nach der
in Abschnitt~\ref{subsec:bewegungsoekologie} beschriebenen von-Mises-Verteilung der
Drehwinkel gewichtet. Woher der Konzentrationsparameter stammt, beschreibt
Abschnitt~\ref{subsec:richtungstreue}.

Der reale Schritt bleibt dabei bewusst Teil der Kandidatenmenge. In der Begriffswelt der
Step-Selection-Funktionen ist er der \emph{genutzte} Schritt unter den \emph{verfügbaren};
ihn auszuschließen würde die Simulation künstlich von der Realität wegdrücken.

\paragraph{Angaben während des Abspielens.} Die Komponente zeigt fortlaufend die
zurückgelegte Strecke, die momentane Geschwindigkeit und den mittleren Abstand zwischen dem
realen und den simulierten Bären. 

\begin{wrapfigure}{r}{0.48\textwidth}
    \centering
    \includegraphics[width=0.46\textwidth]{AnimationsDashboard.png}
    \caption{Animations-Dashboard mit Live-Metriken.}
    \label{fig:animations_dashboard}
\end{wrapfigure}

Die letzte Angabe übersetzt den visuellen Vergleich in eine
ablesbare Zahl (vergleiche die Live-Metriken in Abbildung~\ref{fig:animations_dashboard}). Sie wächst im Verlauf eines Tages systematisch an, und das ist kein Mangel:
Die Simulation wird zu keinem Zeitpunkt auf die reale Route zurückgeholt, sodass sich
Abweichungen zwangsläufig aufsummieren.

Die Animation ist nicht Teil der Hauptkarte, sondern eine eigene Seite mit einer eigenen
Kartenkomponente. Dafür sprechen drei Gründe.

Erstens unterscheidet sich der Datenumfang grundlegend. Die Hauptkarte zeigt Übersichten
über Jahre hinweg und muss die Punktzahl dafür ausdünnen. Die Animation zeigt einen einzelnen
Tag, diesen aber in voller Auflösung und ohne Ausdünnung, weil jede Ortung einen Zeitpunkt
der Bewegung markiert. Beides in einer Komponente zu vereinen hieße, für jede Ansicht den
jeweils falschen Kompromiss zu schließen.

Zweitens unterscheidet sich die Bedienung. Die Animation benötigt eine Zeitleiste,
Abspielsteuerung, Geschwindigkeitsstufen und eine Auswahl des Zeitfensters. Diese Elemente
auf der Hauptkarte unterzubringen würde deren Bedienoberfläche überfrachten, ohne den
übrigen Datenebenen zu nützen.

Drittens unterscheidet sich die Blickrichtung. Auf der Hauptkarte liegen mehrere Bestände
übereinander und werden räumlich verglichen; dafür ist die senkrechte Draufsicht geeignet.
Bei einem einzelnen Bewegungsverlauf ist das Gelände dagegen Teil der Fragestellung, weil
sich erst daran beurteilen lässt, ob ein Pfad einem Tal folgt oder einen Hang quert. Die
Animationsseite ist deshalb auf das Luftbild mit zuschaltbarer Geländedarstellung ausgelegt.

Geteilt werden dennoch die tragenden Bausteine: dieselbe Kartenbibliothek, dieselben
Luftbild- und Höhenquellen, dasselbe Verhalten bei nicht erreichbaren Höhendaten und dieselbe
Tabellendarstellung der zugrunde liegenden Werte unterhalb der Karte. Die Seite zeigt
außerdem an, ob die Daten aus der Projektdatenbank oder aus der lokalen Rückfallebene
stammen - eine Angabe, die für die Nachvollziehbarkeit einer Auswertung wesentlich ist.

```latex
\section{Übertragung des Bewegungsmodells auf die japanischen Kameradaten}
\label{sec:uebertragung}

Das in den vorhergehenden Abschnitten beschriebene Bewegungsmodell basiert auf
Trajektoriendaten aus den rumänischen Karpaten. Für das japanische
Untersuchungsgebiet liegt dagegen keine vergleichbare Menge kontinuierlicher
Bewegungstrajektorien vor. Als Beobachtungsgrundlage stehen dort ausschließlich
punktuelle Detektionen an Kamerastandorten zur Verfügung. Eine direkte Übertragung
der ursprünglichen Modellparameter ist daher nicht ausreichend, da insbesondere
die räumlichen Eigenschaften des Untersuchungsgebiets sowie die Art der verfügbaren
Beobachtungen voneinander abweichen.

Die Übertragung verfolgt deshalb einen modularen Ansatz. Die aus den Karpaten
abgeleiteten Bewegungscharakteristika werden als Grundlage für die Erzeugung
möglicher Schritte verwendet. Die Auswahl zwischen diesen Kandidatenschritten wird
um eine probabilistische Bewertungsfunktion erweitert. Zusätzlich werden
Richtungstreue und Schrittlänge empirisch aus den vorhandenen Trajektoriendaten
abgeleitet. Für das japanische Untersuchungsgebiet wird die Bewegungsbewertung
außerdem um räumliche Informationen aus OpenStreetMap ergänzt. Dadurch werden sowohl
die intrinsischen Bewegungseigenschaften der Tiere als auch die Eigenschaften der
umgebenden Landschaft berücksichtigt.

\subsection{Probabilistische Auswahl der Kandidatenschritte}
\label{subsec:neuimplementierung}

Das ursprüngliche Bewegungsmodell erzeugt für jeden Zeitschritt eine Menge möglicher
Kandidatenschritte. Für die Konstruktion einer vollständigen simulierten Trajektorie
ist zusätzlich eine Regel erforderlich, nach der aus dieser Menge ein konkreter
Schritt ausgewählt wird. Diese Auswahl wird für die Übertragung als probabilistische
Step-Selection-Funktion formuliert.

Für einen Kandidatenschritt $i$ wird zunächst eine Bewertung $s_i$ bestimmt. Sie
beschreibt, wie geeignet der Kandidat unter Berücksichtigung der aktuell
verwendeten Bewegungs- und Geländemerkmale ist. Die Auswahlwahrscheinlichkeit ergibt
sich anschließend aus

\begin{equation}
P_i =
\frac{\exp(\beta s_i)}
{\sum_{j=1}^{n}\exp(\beta s_j)},
\label{eq:step_selection}
\end{equation}

wobei $n$ die Anzahl der verfügbaren Kandidatenschritte und $\beta$ einen
Skalierungsparameter der Bewertungsfunktion bezeichnet.

Die exponentielle Form bewirkt, dass auch Kandidaten mit einer geringeren Bewertung
eine von null verschiedene Auswahlwahrscheinlichkeit besitzen. Gleichzeitig werden
größere Unterschiede zwischen den Bewertungen durch die Exponentialfunktion
verstärkt. Im Gegensatz zu einer deterministischen Auswahl des jeweils am höchsten
bewerteten Kandidaten bleibt damit der stochastische Charakter der Bewegung erhalten.

Die Bewertung kann aus mehreren Komponenten zusammengesetzt werden. Für die
vorliegende Übertragung werden insbesondere die Abweichung von der bisherigen
Bewegungsrichtung sowie die räumliche Eignung des zurückzulegenden Weges
berücksichtigt. Die einzelnen Beiträge können unabhängig voneinander aktiviert
werden. Dadurch lässt sich zwischen einer ausschließlich bewegungsbasierten und
einer zusätzlich geländebasierten Simulation unterscheiden.

Die Kandidatenrichtungen werden geschichtet erzeugt. Der Vollkreis wird hierzu in
acht Sektoren mit einer Breite von jeweils $45^\circ$ unterteilt. Aus jedem Sektor
wird genau eine Richtung zufällig bestimmt. Die Streuung innerhalb eines Sektors
beträgt $\pm22{,}5^\circ$. Auf diese Weise wird der gesamte Winkelbereich abgedeckt,
während gleichzeitig vermieden wird, dass mehrere Kandidaten zufällig in einem
kleinen Winkelbereich konzentriert sind.

Die Auswahlwahrscheinlichkeit eines Kandidaten wird damit nicht allein durch seine
Richtung bestimmt. Vielmehr entsteht für jeden Zeitschritt ein diskreter
Möglichkeitsraum, in dem die verschiedenen Kandidaten anhand ihrer Eigenschaften
unterschiedlich wahrscheinlich ausgewählt werden.

\subsection{Empirische Modellierung der Richtungstreue}
\label{subsec:richtungstreue}

Ein wesentliches Merkmal tierischer Bewegungen ist die Abhängigkeit der aktuellen
Bewegungsrichtung von der Richtung des unmittelbar vorausgegangenen Schrittes.
Diese Richtungstreue wird im Modell über eine von-Mises-Verteilung beschrieben.
Der zugehörige Konzentrationsparameter $\kappa$ bestimmt dabei, wie stark die
Auswahl auf bestimmte Richtungsänderungen konzentriert ist.

Anstatt $\kappa$ als festen Wert vorzugeben, wird die Richtungstreue aus den
Karpaten-Trajektorien empirisch abgeleitet. Grundlage sind 46\,161 Paare
aufeinanderfolgender Schritte von 13 Individuen. Für jedes Schrittpaar wird der
Drehwinkel $\Delta\theta$ zwischen den beiden Bewegungsrichtungen bestimmt und
anschließend dessen Kosinus betrachtet. Als Maß für die mittlere Richtungstreue
ergibt sich

\begin{equation}
\rho =
\frac{1}{n}
\sum_{i=1}^{n}
\cos(\Delta\theta_i).
\label{eq:rho_empirical}
\end{equation}

Ein Wert von $\rho=1$ entspricht dabei einer vollständigen Beibehaltung der
Bewegungsrichtung, während $\rho=0$ keine gerichtete Abhängigkeit beschreibt.
Negative Werte weisen auf eine Tendenz zu einer Richtungsumkehr hin.

Die Auswertung zeigt, dass die Richtungstreue deutlich von der Schrittlänge
abhängt. Die entsprechenden Werte sind in Tabelle~\ref{tab:richtungstreue}
zusammengefasst.

\begin{table}[htbp]
\centering
\caption{Empirische Richtungstreue der Karpatenbären in Abhängigkeit von der
Schrittlänge}
\label{tab:richtungstreue}
\small
\begin{tabular}{l r r r}
\toprule
\textbf{Schrittlänge} &
\textbf{Anzahl} &
\textbf{mittlerer Kosinus} &
\textbf{Median des Drehwinkels} \\
\midrule
unter 25\,m      & 24.058 & $-0{,}202$ & $114^\circ$ \\
25 bis 50\,m     &  5.876 & $-0{,}122$ & $106^\circ$ \\
50 bis 100\,m    &  4.545 & $+0{,}106$ & $75^\circ$ \\
100 bis 200\,m   &  4.529 & $+0{,}251$ & $62^\circ$ \\
200 bis 500\,m   &  4.634 & $+0{,}414$ & $47^\circ$ \\
über 500\,m      &  2.519 & $+0{,}588$ & $36^\circ$ \\
\bottomrule
\end{tabular}
\end{table}

Über alle Schritte ergibt sich ein Mittelwert von $\rho=-0{,}012$. Dieser
aggregierte Wert würde zunächst auf eine nahezu richtungsunabhängige Bewegung
hindeuten. Die Aufteilung nach Schrittlänge zeigt jedoch einen systematischen
Zusammenhang: Bei kurzen Schritten treten überwiegend große Richtungsänderungen
auf, während längere Schritte eine deutlich stärkere Beibehaltung der bisherigen
Richtung aufweisen.

Eine über alle Schrittlängen konstante Richtungstreue würde diese unterschiedlichen
Bewegungszustände nur unzureichend abbilden. Daher wird $\rho$ in der Simulation
abhängig von der erwarteten Schrittlänge bestimmt. Zwischen den empirischen
Stützstellen wird auf logarithmischer Schrittlängenachse interpoliert. Außerhalb
des empirisch erfassten Bereichs wird der jeweils nächstgelegene Randwert
verwendet, sodass keine Extrapolation erforderlich ist.

Der Zusammenhang zwischen der mittleren Richtungstreue $\rho$ und dem
Konzentrationsparameter $\kappa$ der von-Mises-Verteilung wird über

\begin{equation}
\rho = A(\kappa)
=
\frac{I_1(\kappa)}{I_0(\kappa)}
\label{eq:rho}
\end{equation}

beschrieben. Dabei bezeichnen $I_0$ und $I_1$ modifizierte Besselfunktionen
erster Art. Für die Simulation wird dieser Zusammenhang invertiert, um aus dem
empirisch bestimmten $\rho$ den entsprechenden Wert von $\kappa$ zu erhalten.
Hierzu wird die Näherung nach Fisher verwendet.

Die Richtungsbewertung eines Kandidaten ergibt sich damit aus

\begin{equation}
w_{\mathrm{dir}}
=
\exp\left(
\kappa \cos(\Delta\theta)
\right).
\label{eq:direction_weight}
\end{equation}

Bei positivem $\kappa$ werden Richtungen bevorzugt, die der bisherigen
Bewegungsrichtung entsprechen. Ein negatives $\kappa$ führt dagegen zu einer
Bevorzugung größerer Richtungsänderungen und kann somit das bei kurzen Schritten
beobachtete Hin-und-Her-Verhalten abbilden.

Die empirische Kalibrierung ist allerdings mit einer Unsicherheit behaftet. Der
Median der betrachteten Schrittlängen beträgt 22{,}7\,m und liegt damit in einer
Größenordnung, in der sich die Positionsungenauigkeit der GPS-Daten auf die
berechneten Richtungsänderungen auswirken kann. Insbesondere die negativen Werte
bei kurzen Schritten sind daher nicht ausschließlich als tatsächliches
Bewegungsverhalten zu interpretieren. Dieser Aspekt wird bei der Bewertung des
Modells in Abschnitt~\ref{sec:limitationen} berücksichtigt.

\subsection{Geländebasierte Bewertung der Kandidatenschritte}
\label{subsec:gelaende}

Die bewegungsbasierte Auswahl berücksichtigt die bisherige Bewegungsrichtung,
enthält jedoch zunächst keine Information über die räumliche Umgebung. Für die
Übertragung auf das japanische Untersuchungsgebiet wird deshalb eine
geländebasierte Bewertung ergänzt.

Die Grundlage hierfür bilden OpenStreetMap-Daten, die für das Untersuchungsgebiet
in ein regelmäßiges Raster überführt werden. Jeder Rasterzelle wird anhand ihrer
Geländemerkmale eine Attraktivitätsbewertung zugeordnet. Potenziell geeignete
Lebensräume, beispielsweise Wald, Gebüsch oder Schutzgebiete, tragen positiv zur
Bewertung bei. Anthropogene Nutzungsformen sowie bestimmte Verkehrswege werden
dagegen negativ bewertet. Straßen werden dabei entsprechend ihrer
Klassifizierung unterschiedlich berücksichtigt.

Für die Bewertung werden nicht ausschließlich die Werte einzelner
OpenStreetMap-Tags betrachtet. Entscheidend ist vielmehr die Kombination aus
Schlüssel und Wert eines Merkmals. Dies ist insbesondere bei Merkmalen wie
\texttt{residential} relevant, da ein Wohngebiet und eine als
\texttt{highway=residential} klassifizierte Wohnstraße unterschiedliche
räumliche Bedeutungen besitzen.

Die Geländebewertung eines Kandidatenschrittes wird nicht ausschließlich am
Endpunkt bestimmt. Stattdessen werden mehrere Positionen entlang des möglichen
Weges betrachtet. Dadurch fließen auch Geländeeigenschaften zwischen Ausgangs-
und Zielpunkt in die Bewertung ein.

Diese Vorgehensweise ist insbesondere aufgrund des Verhältnisses zwischen
Schrittlänge und Rasterauflösung relevant. Die typischen Bewegungsschritte sind
teilweise kürzer als die verwendete Rasterweite von ungefähr 150\,m. Eine
ausschließliche Bewertung des Endpunkts würde deshalb bei mehreren Kandidaten
häufig auf dieselbe Rasterzelle führen und den Einfluss des Geländes reduzieren.

Die Bewertung entlang des Weges berücksichtigt außerdem mögliche Hindernisse
zwischen dem aktuellen Standort und einem attraktiven Zielgebiet. Ein attraktives
Waldgebiet hinter einer Siedlung soll beispielsweise nicht automatisch zu einer
Bevorzugung der direkten Richtung führen, wenn der Weg dorthin durch eine
ungünstig bewertete Fläche verläuft.

\subsection{Lokale Bereitstellung der Geländedaten}
\label{subsec:terrain_cache}

Die für die Geländebewertung benötigten OpenStreetMap-Merkmale werden über die
Overpass-Schnittstelle bezogen. Eine Abfrage für jeden Kandidatenschritt würde
jedoch zu einer großen Anzahl einzelner Netzwerkanfragen führen und die
Simulation von der Verfügbarkeit und Antwortzeit eines externen Dienstes
abhängig machen.

Daher wird das Untersuchungsgebiet vor der Simulation vollständig abgefragt und
lokal als Raster vorgehalten. Während eines Simulationslaufs werden die
Geländebewertungen innerhalb dieses vorbereiteten Gebiets ausschließlich aus
dem lokalen Datenspeicher gelesen. Die eigentliche Bewegungssimulation ist damit
von der externen Datenbeschaffung entkoppelt.

Diese Trennung besitzt neben der Laufzeiteffizienz auch eine methodische
Bedeutung. Die Eingangsdaten der Simulation bleiben während eines Laufs konstant
und können für wiederholte Experimente erneut verwendet werden. Dadurch wird
vermieden, dass unterschiedliche Ergebnisse allein durch zeitlich veränderte
Antworten des externen Kartendienstes entstehen.

Für Positionen außerhalb des vorbereiteten Gebiets ist in der Implementierung
eine externe Abfrage als Rückfallebene vorgesehen. Für die eigentliche
quantitative Auswertung muss jedoch sichergestellt werden, dass das vorbereitete
Gebiet sämtliche während der Simulation erreichbaren Positionen umfasst. Andernfalls
kann ein Wechsel zwischen lokaler und externer Datenquelle zu unterschiedlichen
Geländegrundlagen innerhalb eines Simulationslaufs führen.

\subsection{Clusterbasierte Richtungswahl}
\label{subsec:cluster}

Die lokale Geländebewertung beschreibt die Eigenschaften einzelner Kandidatenwege.
Zusätzlich soll berücksichtigt werden können, ob sich in einer bestimmten Richtung
ein größeres zusammenhängendes Gebiet geeigneten Habitats befindet. Hierzu wird
das Prinzip der gewichteten Schwerpunktbildung aus Abschnitt~\ref{subsec:schwerpunkt}
auf das Geländeraster übertragen.

Die zugrunde liegende Berechnung basiert auf

\begin{equation}
P =
\frac{\sum_i w_i p_i}
{\sum_i w_i},
\label{eq:schwerpunkt}
\end{equation}

wobei $p_i$ die Position einer Rasterzelle und $w_i$ deren Gewicht bezeichnet.
Der berechnete Punkt stellt den gewichteten räumlichen Schwerpunkt der
berücksichtigten Zellen dar.

Vor der Schwerpunktbildung werden die attraktiven Rasterzellen räumlich getrennt.
Hierzu werden Zellen oberhalb einer definierten Attraktivitätsschwelle anhand
ihrer räumlichen Nachbarschaft zu zusammenhängenden Gebieten gruppiert. Die
Implementierung verwendet hierfür eine Zusammenhangskomponentensuche mit
8er-Nachbarschaft.

Diese Vorgehensweise stellt eine Anpassung an die Struktur der verfügbaren
Geländedaten dar. Die zugrunde liegende Methode aus der Referenz verwendet für
die Trennung der relevanten Bereiche einen gelernten Merkmalsraum. Ein solcher
Merkmalsraum liegt für das hier verwendete regelmäßige Geländeraster nicht vor.
Die räumliche Nachbarschaft der Rasterzellen ist hingegen unmittelbar gegeben,
sodass die Zusammenhangskomponentensuche die räumliche Trennung direkt auf dem
Raster durchführen kann.

Übertragen wird somit das Prinzip, zunächst getrennte Bereiche zu identifizieren
und anschließend innerhalb dieser Bereiche einen gewichteten Schwerpunkt zu
bestimmen. Die konkrete Methode der Gebietsbildung wird dagegen an die
Struktur der vorliegenden Daten angepasst.

\paragraph{Positive Gewichte für die Schwerpunktbildung.}

Die Geländebewertungen enthalten sowohl positive als auch negative Beiträge.
Eine direkte Verwendung dieser Bewertungen als Gewichte in
Gleichung~\eqref{eq:schwerpunkt} wäre deshalb nicht sinnvoll, da sich positive
und negative Beiträge gegenseitig aufheben könnten. Insbesondere könnte der
Nenner der Gleichung gegen null gehen und dadurch zu instabilen
Schwerpunktpositionen führen.

Für die Schwerpunktbildung werden deshalb ausschließlich Rasterzellen
oberhalb der Attraktivitätsschwelle berücksichtigt. Ihr Gewicht wird relativ
zu dieser Schwelle bestimmt:

\begin{equation}
w_i =
s_i-s_{\mathrm{threshold}}+0{,}5.
\label{eq:cluster_weight}
\end{equation}

Damit bleiben die verwendeten Gewichte positiv. Die negativen Geländebeiträge
werden weiterhin bei der Bewertung einzelner Kandidatenschritte berücksichtigt,
gehen jedoch nicht als negative Masse in die Schwerpunktbildung ein.

Für jedes zusammenhängende Gebiet wird anschließend der gewichtete Schwerpunkt
bestimmt. Die Richtung vom aktuellen Simulationsstandort zu diesem Schwerpunkt
wird als zusätzlicher Kandidat in die Schrittauswahl aufgenommen.

\paragraph{Berücksichtigung der räumlichen Masse.}

Die Richtung zu einem zusammenhängenden Gebiet stellt zunächst lediglich einen
zusätzlichen Kandidaten dar. Damit die Information über die Größe und
Attraktivität des Gebiets tatsächlich in die Auswahl eingeht, wird die
Richtung zusätzlich anhand der Masse des jeweiligen Gebiets gewichtet.

Der hierfür verwendete Bonus steigt mit der Clustergröße beziehungsweise
Clusterbewertung und wird durch eine Tangens-hyperbolicus-Funktion begrenzt.
Damit kann ein größeres attraktives Gebiet die Wahrscheinlichkeit seiner
Richtung erhöhen, ohne die übrigen Kandidaten vollständig aus dem
Auswahlprozess zu verdrängen.

Die freien Kandidatenrichtungen bleiben gleichzeitig Bestandteil der
Kandidatenmenge. Dadurch wird vermieden, dass sich die Simulation ausschließlich
auf die identifizierten Habitatgebiete beschränkt. Die Bewegungsentscheidung
bleibt somit eine Kombination aus lokaler Bewegungsdynamik, Geländequalität und
räumlich aggregierter Habitatinformation.

\subsection{Zusammenführung der Modellkomponenten}
\label{subsec:integration_uebertragung}

Die übertragene Bewegungslogik verbindet die beschriebenen Komponenten zu einer
gemeinsamen probabilistischen Schrittauswahl. Für jeden Simulationsschritt wird
zunächst anhand des jeweiligen Zeitintervalls die erwartete Schrittlänge
bestimmt. Auf dieser Grundlage werden Kandidatenschritte mit unterschiedlichen
Richtungen erzeugt.

Die Richtungsbewertung berücksichtigt anschließend die Abweichung von der
bisherigen Bewegungsrichtung. Der hierfür verwendete Konzentrationsparameter
$\kappa$ wird aus der empirisch beobachteten Richtungstreue und der erwarteten
Schrittlänge bestimmt. Sofern Geländedaten verfügbar sind, wird zusätzlich die
Geländequalität entlang des Kandidatenschrittes bewertet.

Werden in der Umgebung zusammenhängende attraktive Gebiete identifiziert, werden
deren gewichtete Schwerpunkte als zusätzliche Richtungsalternativen in die
Kandidatenmenge aufgenommen. Ihre räumliche Masse beeinflusst die Bewertung über
den begrenzten Clusterbonus.

Die resultierende Gesamtbewertung $s_i$ eines Kandidaten wird schließlich über
Gleichung~\eqref{eq:step_selection} in eine Auswahlwahrscheinlichkeit überführt.
Die Bewegung des simulierten Individuums ergibt sich somit aus einer zufälligen
Ziehung innerhalb eines durch Bewegungsrichtung, Schrittlänge und
Geländeeigenschaften eingeschränkten Möglichkeitsraums.

Ein wesentlicher Bestandteil der Übertragung ist dabei die Trennung zwischen
empirisch abgeleiteten und modellseitig eingeführten Komponenten. Die
Richtungstreue wird aus den vorhandenen Karpaten-Trajektorien bestimmt, während
die Geländegewichtung und die clusterbasierte Richtungswahl aus den Eigenschaften
des japanischen Untersuchungsgebiets abgeleitet werden. Damit wird nicht
behauptet, dass die japanischen Kameradaten selbst die zugrunde liegenden
Bewegungsparameter liefern. Vielmehr wird ein aus den verfügbaren
Trajektoriendaten abgeleitetes Bewegungsmodell unter Verwendung der im
japanischen Untersuchungsgebiet verfügbaren räumlichen Informationen
angewendet.

Die resultierenden simulierten Trajektorien bilden anschließend die Grundlage
für die zeitliche Gegenüberstellung mit den realen Kameradetektionen. Da aus
den punktuellen Kameradetektionen keine vollständigen realen Trajektorien
rekonstruiert werden können, erfolgt die eigentliche quantitative Bewertung
nicht über einen direkten räumlichen Trajektorienvergleich, sondern über die
in Abschnitt~\ref{sec:tagesrhythmus} beschriebene zeitliche Analyse.
```

\section{Statistischer Tagesrhythmus-Vergleich}
\label{sec:tagesrhythmus}

Der vorangegangene Abschnitt hat das Bewegungsmodell auf die japanische Landschaft
übertragen. Offen bleibt, wie sich das Ergebnis überhaupt prüfen lässt. Ein räumlicher
Pfadvergleich scheidet aus, denn aus punktuellen Kameradetektionen lässt sich keine
Trajektorie rekonstruieren (Abschnitt~\ref{subsec:kamerafallen}). Der Vergleich wird deshalb
auf die zeitliche Ebene verlagert: Simulierte und reale Daten besitzen beide eine
Tageszeitverteilung, und diese lassen sich gegenüberstellen.

\subsection{Zirkuläre Kerndichteschätzung und Überlappungskoeffizient}
\label{subsec:kde-umsetzung}

\paragraph{Reale Auslösezeiten.} Die realen Zeitpunkte stammen aus der vereinheitlichten
Bildtabelle. Ausgeschlossen werden Aufnahmen zweier Gerätekennungen, die dem Muster der
übrigen sechsstelligen Kennungen nicht folgen und mit hoher Wahrscheinlichkeit aus einer
fehlgeschlagenen Kennungserkennung beim E-Mail-Abruf stammen. Die verbleibenden Zeitstempel
werden in Dezimalstunden umgerechnet.

\paragraph{Zusammenfassung über alle Kameras.} Die Auswertung erfolgt über alle Kameras
gemeinsam und nicht je Kamera. Das ist eine Notwendigkeit, keine Vereinfachung: Die realen
Daten umfassen nur wenige Dutzend verwertbare Zeitstempel, verteilt auf rund ein bis zwei
Dutzend Kameras. Für eine Verteilung \emph{je Kamera} ist das bei weitem zu wenig.
Zusammengefasst reicht die Stichprobe für eine Kerndichteschätzung, die aber weiterhin mit
Vorsicht zu lesen ist; Abschnitt~\ref{sec:limitationen} ordnet das ein.

\paragraph{Umsetzung der zirkulären Schätzung.} Umgesetzt ist die in
Abschnitt~\ref{subsec:kamerafallen} beschriebene Kerndichteschätzung auf dem 24-Stunden-Kreis.
Der Ringschluss wird dadurch erreicht, dass die Beobachtungen zusätzlich um $-24$ und $+24$
Stunden versetzt in die Schätzung eingehen. Damit liegen 23:30~Uhr und 00:30~Uhr korrekt
nahe beieinander und nicht an entgegengesetzten Enden einer Geraden.

Beide Dichten werden auf einem gemeinsamen Gitter von 481~Stützstellen über 24~Stunden
ausgewertet und so normiert, dass ihr Integral eins ergibt. Der Überlappungskoeffizient nach
Gleichung~\eqref{eq:delta} wird anschließend als Summe der punktweisen Minima über dieses
Gitter genähert.

\paragraph{Abgrenzung zur Referenzumsetzung.} Das Verfahren folgt dem Prinzip von Ridout und
Linkie \cite{RidoutLinkie2009}, verwendet aber einen umwickelten Gauß-Kern statt des
von-Mises-Kerns des dort beschriebenen R-Pakets. Diese Abweichung ist zu benennen: Die
Zahlenwerte sind dadurch nicht unmittelbar mit Ergebnissen vergleichbar, die mit jenem Paket
erzeugt wurden.

\subsection{Monte-Carlo-Simulation virtueller Kamera-Detektionen}
\label{subsec:montecarlo}

Der realen Verteilung muss eine simulierte gegenübergestellt werden. Dafür genügt kein
einzelner Simulationslauf, denn ein Pfad ist eine Zufallsziehung. Erzeugt werden deshalb
viele unabhängige Läufe.

\paragraph{Aufbau eines Laufs.} Je Lauf wird ein Bär an einem zufällig gewählten
Kamerastandort gesetzt und über 48~Stunden bewegt. Geschlecht, Altersklasse und Jahresphase
werden zufällig gezogen. Das ist eine bewusste Entscheidung: Es ist nicht bekannt, welches
Tier welche Kamera ausgelöst hat, und die Zufallsziehung mittelt über die unbekannte
Zusammensetzung der Population, statt eine anzunehmen. Ebenso wird die Startstunde
gleichverteilt über den Tag gezogen, damit die Startzeit die spätere Tagesverteilung nicht
vorprägt.

\paragraph{Die virtuelle Auslösung.} Nähert sich der simulierte Bär einer beliebigen Kamera
auf weniger als den Erfassungsradius, zählt das als virtuelle Auslösung und die zugehörige
Tageszeit wird aufgenommen. Je Zeitschritt wird höchstens eine Auslösung gezählt, auch wenn
mehrere Kameras in Reichweite liegen.

Der Startpunkt selbst zählt nicht mit. Er ist definitionsgemäß eine Kamera, sodass andernfalls
jeder Lauf trivialerweise sofort \enquote{erfasst} würde und die Verteilung durch die
gleichverteilt gezogene Startstunde verfälscht wäre.

\paragraph{Der Beispiellauf in der Oberfläche.} Für die Darstellung auf der Karte wird einer
der Läufe mitgeschnitten. Entscheidend ist die Auswahlregel: Es ist \emph{nicht} der Lauf mit
den meisten Treffern, was eine Auswahl zugunsten der schöneren Vorführung wäre, sondern
schlicht der erste Lauf mit mindestens einem Treffer. Gesucht wird nur innerhalb der ersten
15~Läufe, damit die Darstellung nicht zufällig einen völlig leeren Lauf zeigt; bleiben auch
diese ohne Treffer, wird der letzte davon gezeigt und in der Oberfläche als Lauf ohne Treffer
gekennzeichnet. Die Oberfläche bezeichnet ihn durchgehend als einen zufälligen Lauf und nicht
als bestes Ergebnis.

\subsection{Unsicherheitsquantifizierung mittels Bootstrap}
\label{subsec:bootstrap-umsetzung}

Ein einzelner Zahlenwert für~$\Delta$ suggeriert eine Genauigkeit, die bei einer Stichprobe
von wenigen Dutzend realen Zeitpunkten nicht gerechtfertigt wäre. Angegeben wird deshalb ein
Intervall.

Verwendet wird das in Abschnitt~\ref{subsec:kamerafallen} beschriebene Bootstrap-Verfahren:
Aus beiden Stichproben wird jeweils mit Zurücklegen eine neue Stichprobe desselben Umfangs
gezogen, daraus werden beide Dichten neu geschätzt und~$\Delta$ neu berechnet. Das geschieht
1000~mal. Aus der Verteilung dieser 1000~Werte ergeben sich das 2{,}5\%- und das
97{,}5\%-Quantil als Grenzen eines 95\%-Intervalls.

Was das Intervall beantwortet, lässt sich einfach sagen: Wie stark würde sich das Ergebnis
ändern, wenn die Stichprobe geringfügig anders ausgefallen wäre? Ein breites Intervall ist
deshalb kein Mangel der Auswertung, sondern eine zutreffende Aussage über die Datenlage.

Die Oberfläche zeigt zusätzlich die Verteilung der 1000~Bootstrap-Werte als Histogramm mit
eingezeichnetem Punktschätzer und Intervallgrenzen. Das macht unmittelbar sichtbar, wie breit
die Streuung tatsächlich ist, statt sie auf zwei Zahlen zu reduzieren.

\paragraph{Reproduzierbarkeit und gemeinsame Berechnung.} Sämtliche Zufallsziehungen gehen
von einem einstellbaren Startwert aus; ein Lauf ist damit exakt wiederholbar
(Anforderung~N5). Die vollständige Berechnung ist zudem an genau einer Stelle umgesetzt und
wird sowohl vom Kommandozeilenskript als auch von der Oberfläche aufgerufen. Damit ist
ausgeschlossen, dass zwei Umsetzungen nebeneinander bestehen, deren Ergebnisse auseinanderlaufen.

% ====================================================================
\chapter{Diskussion}
\label{chap:diskussion}

Dieses Kapitel bewertet das Ergebnis. Es gleicht zunächst das entwickelte System gegen die in
Abschnitt~\ref{sec:anforderungen} aufgestellten Anforderungen ab, ordnet anschließend das
ermittelte Übereinstimmungsmaß ein und benennt zuletzt die Grenzen der Arbeit.

\section{Erfüllung der Anforderungen und Bewertung des Systems}
\label{sec:anforderungserfuellung}

\subsection{Abgleich gegen die Anforderungen}
\label{subsec:abgleich}

Tabelle~\ref{tab:anforderungserfuellung} stellt den Erfüllungsgrad je Anforderung dem
Abschnitt gegenüber, in dem die Umsetzung beschrieben ist. Von den 17~Anforderungen sind 15
vollständig erfüllt; die beiden übrigen werden im Anschluss einzeln begründet.

\begin{table}[htbp]
\centering
\caption{Erfüllung der Anforderungen aus Abschnitt~\ref{sec:anforderungen}}
\label{tab:anforderungserfuellung}
\small
\begin{tabular}{L{1.1cm} L{6.6cm} L{2.4cm} L{3.0cm}}
\toprule
\textbf{Nr.} & \textbf{Anforderung} & \textbf{Erfüllung} & \textbf{Abschnitt} \\
\midrule
F1 & Gemeinsame Karte & erfüllt & \ref{subsec:aufbau-karte} \\
F2 & Filterung & erfüllt & \ref{sec:interaktion} \\
F3 & Eigene Marker und Flächen & erfüllt & \ref{sec:interaktion} \\
F4 & Zeitlicher Verlauf & erfüllt & \ref{sec:animation} \\
F5 & Gegenüberstellung von Modell und Beobachtung & erfüllt & \ref{subsec:synchron} \\
F6 & Erzeugung zusammenhängender Pfade & erfüllt & \ref{subsec:neuimplementierung} \\
F7 & Statistischer Vergleich & erfüllt & \ref{sec:tagesrhythmus} \\
F8 & Einsicht in die Rohdaten & erfüllt & \ref{sec:interaktion} \\
F9 & Sichtbarkeit neuer Aufnahmen & erfüllt & \ref{subsec:anbindung} \\
\addlinespace
N1 & Darstellungsleistung & erfüllt & \ref{subsec:aufbau-karte} \\
N2 & Keine schwergewichtigen Abhängigkeiten & erfüllt & \ref{subsec:anbindung} \\
N3 & Einheitliche Datenhaltung & erfüllt & \ref{sec:datenhaltung} \\
N4 & Nachvollziehbare Fehler & erfüllt & \ref{sec:datenhaltung} \\
N5 & Reproduzierbarkeit & erfüllt & \ref{subsec:bootstrap-umsetzung} \\
N6 & Übertragbare Konfiguration & teilweise & \ref{sec:datenhaltung} \\
N7 & Einheitlicher Bildbestand & erfüllt & \ref{sec:datenhaltung} \\
N8 & Eindeutigkeit der Ortungen & offen & \ref{subsec:technische-grenzen} \\
\bottomrule
\end{tabular}
\end{table}

\paragraph{N6 - teilweise erfüllt.} Die Zugangsdaten werden aus Umgebungsvariablen gelesen,
sodass sich das System ohne Eingriff in den Quelltext auf einem anderen Rechner betreiben
lässt. Die bisherigen Werte sind jedoch als Rückfallebene noch im Quelltext hinterlegt. Das
war während der Entwicklung sinnvoll, weil bestehende Arbeitsplätze unverändert weiterliefen,
muss vor einem produktiven Einsatz aber entfallen. Hinzu kommt, dass das Zugangswort in der
Versionsgeschichte des Projekts enthalten ist und deshalb unabhängig von dieser Änderung zu
wechseln ist.

\paragraph{Anmerkung zu N1.} Die Darstellung bleibt auch dann flüssig bedienbar, wenn
mehrere Datenebenen mit zusammen mehreren zehntausend Objekten gleichzeitig sichtbar sind;
die Anforderung gilt damit als erfüllt. Diese Beurteilung beruht jedoch auf der Erprobung im
Betrieb und nicht auf einer systematischen Messung mit festgelegtem Messaufbau. Ein Vergleich
beider Technologien unter gleichen Bedingungen wurde nicht durchgeführt, da dazu die
abgelöste Darstellung für denselben Datenumfang hätte nachgebaut werden müssen. Die Aussage
ist entsprechend als praktische Beobachtung zu lesen.

\paragraph{N8 - offen.} Die in Abschnitt~\ref{subsec:datenhaltung} festgestellte
Mehrfachspeicherung der Karpaten-Ortungen wird beim Lesen nicht gefiltert. Die Ursache ist
bekannt und der Eingriff wäre klein; er unterblieb, weil er sich erst spät im Verlauf der
Arbeit zeigte und eine Änderung an der zentralen Datenzugriffsschicht ohne anschließende
vollständige Prüfung aller darauf aufsetzenden Auswertungen nicht vertretbar erschien.
Abschnitt~\ref{subsec:technische-grenzen} ordnet die Auswirkung ein.

\subsection{Bewertung der Entwurfsentscheidungen}
\label{subsec:bewertung-entwurf}

Zwei Entscheidungen prägen das System stärker als alle übrigen und verdienen eine
rückblickende Beurteilung.

\paragraph{Die statische Kartenkomponente.} Die Entscheidung gegen eine bidirektionale
Komponente (Abschnitt~\ref{subsec:architekturentscheidung}) hat sich im Wesentlichen
bewährt. Der weitaus größte Teil der Interaktion betrifft allein die Darstellung und läuft
damit ohne jeden Neuaufbau der Seite ab - schneller, als es mit einem Rückkanal der Fall
wäre, da dieser jedes Mal einen Durchlauf auf der Serverseite ausgelöst hätte.

Der Preis ist an genau einer Stelle spürbar. Das Anlegen neuer Marker und Flächen erfordert
den Umweg über die Adresszeile und ein tatsächliches Neuladen der Seite. Das funktioniert,
ist aber umständlicher als ein Rückgabewert, und es setzt voraus, dass eingebettete
Komponente und umgebende Seite derselben Herkunft zugerechnet werden. Hinter einem
vorgeschalteten Server mit abweichendem Namen entfiele diese Voraussetzung, und sowohl der
Zustandserhalt als auch das Zeichnen würden ausfallen - die Karte selbst bliebe nutzbar.
Dieser Fall ist vor einem Ausrollen zu prüfen.

Rückblickend war die Entscheidung dennoch richtig, weil ihr Hauptargument außerhalb der
Technik liegt: Das System soll an der Hochschule weitergeführt werden, und eine Anwendung
ohne getrennten Übersetzungsschritt ist dafür erheblich zugänglicher.

\paragraph{Die Neuimplementierung der Schrittauswahl.} Die Entscheidung, die Auswahl neu zu
schreiben statt die vorhandene zu erweitern (Abschnitt~\ref{subsec:neuimplementierung}), war
weniger eine Wahl als eine Folge des Befunds: Es gab keine Auswahl, die sich hätte erweitern
lassen. Der Nutzen zeigt sich daran, dass die drei wesentlichen Verbesserungen - empirisch
geschätzte Richtungstreue, Geländegewichtung und lückenlose Richtungsabdeckung - erst
dadurch möglich wurden.

Der Preis ist, dass nun zwei Umsetzungen der Schritterzeugung nebeneinander bestehen: die
der Vorarbeit und die dieser Arbeit. Das ist vertretbar, weil die Vorarbeit unverändert
erhalten bleibt und die Abweichungen dokumentiert sind, erfordert bei künftigen Änderungen
aber Aufmerksamkeit.

\section{Wissenschaftliche Einordnung des \texorpdfstring{$\Delta$}{Delta}-Ergebnisses}
\label{sec:einordnung-delta}

\subsection{Was der Wert aussagt}
\label{subsec:was-delta-aussagt}

Der Überlappungskoeffizient misst, wie stark sich zwei Tagesaktivitätsverteilungen decken.
Er nimmt den Wert~0 an, wenn sich die Verteilungen nicht überschneiden, und~1, wenn sie
identisch sind.

Eine erste Einordnung ergibt sich aus der Betrachtung, was ein beliebiges Modell mindestens
erreichen würde. Eine über den Tag gleichverteilte simulierte Aktivität erzielt gegenüber
einer deutlich nachtaktiven realen Verteilung bereits einen merklichen Überlappungswert,
ohne irgendetwas über das Verhalten der Tiere zu wissen. Ein Wert ist deshalb nicht gegen
null zu bewerten, sondern gegen diesen trivialen Vergleichsmaßstab.

\subsection{Was der Wert nicht aussagt}
\label{subsec:was-delta-nicht-aussagt}

Drei Abgrenzungen sind wesentlich, weil sie häufig überlesen werden.

Erstens ist eine Übereinstimmung der Tagesrhythmen \emph{keine} Bestätigung der räumlichen
Pfade. Verglichen wird ausschließlich die zeitliche Dimension. Zwei Modelle mit völlig
verschiedenem Raumverhalten können denselben Tagesrhythmus erzeugen.

Zweitens stammt der Tagesrhythmus der Simulation nicht aus dem Modellverhalten im
eigentlichen Sinne, sondern aus der hinterlegten Geschwindigkeitstabelle. Diese enthält je
Stunde, Geschlecht, Altersklasse und Jahresphase eine Mediangeschwindigkeit und bestimmt
damit, wann sich ein simuliertes Tier überhaupt bewegt und folglich wann es eine Kamera
auslösen kann. Die Übereinstimmung prüft damit in erster Linie, ob die an rumänischen
Bären gemessenen Tagesrhythmen auf die japanischen Verhältnisse übertragbar sind - eine
sinnvolle Frage, aber eine engere als \enquote{funktioniert das Modell}.

Drittens ist der Vergleich durch die Stichprobe begrenzt. Bei wenigen Dutzend realen
Zeitpunkten ist das Bootstrap-Intervall breit, und ein Punktschätzer ohne dieses Intervall
wäre irreführend. Das Ergebnis ist deshalb als explorativer Befund zu lesen und nicht als
belastbarer Nachweis.

\subsection{Einordnung gegenüber der Literatur}
\label{subsec:literatureinordnung}

Das Verfahren selbst ist in der Kamerafallenforschung etabliert
\cite{RidoutLinkie2009, Rowcliffe2014}. Abweichend von der verbreiteten Umsetzung wird hier
ein umwickelter Gauß-Kern statt eines von-Mises-Kerns verwendet
(Abschnitt~\ref{subsec:kde-umsetzung}), weshalb die Zahlenwerte mit Ergebnissen aus jener
Umsetzung nicht unmittelbar vergleichbar sind.

Eine weitere Besonderheit unterscheidet diese Arbeit von der üblichen Anwendung: Dort werden
zwei \emph{beobachtete} Verteilungen verglichen, etwa zweier Arten am selben Standort. Hier
steht einer beobachteten eine \emph{simulierte} Verteilung gegenüber. Der Koeffizient misst
damit nicht die Überschneidung zweier Tiergruppen, sondern die Übereinstimmung zwischen
Modell und Beobachtung. Das ist eine zulässige, aber eigenständige Verwendung, die bei der
Interpretation mitzudenken ist.

\section{Limitationen}
\label{sec:limitationen}

\subsection{Datenbasis}
\label{subsec:limit-daten}

\paragraph{Stichprobenumfang der Kameradaten.} Die schwerwiegendste Einschränkung ist der
Umfang der realen Auslösezeiten. Er reicht für eine zusammengefasste Auswertung über alle
Kameras, nicht aber für eine Betrachtung einzelner Standorte - und gerade die wäre
interessant, da sich Kameras in Waldlage und in Siedlungsnähe im Tagesrhythmus unterscheiden
dürften. Die Literatur zur Schätzung von Aktivitätskurven weist darauf hin, dass der
Stichprobenumfang die Genauigkeit solcher Schätzungen maßgeblich bestimmt
\cite{Rowcliffe2014}.

\paragraph{Aufstellung der Kameras.} Die Kameras stehen nicht zufällig im Gelände, sondern
dort, wo Konflikte mit Menschen erwartet werden, also in Siedlungsnähe und an bekannten
Wechseln. Für die vorliegende Auswertung ist das weniger gravierend als für
Dichteschätzungen, die eine zufällige Aufstellung ausdrücklich voraussetzen
\cite{Rowcliffe2008}. Dennoch bleibt die Einschränkung bestehen: Gemessen wird der
Tagesrhythmus an Orten menschlicher Nutzung, der sich von dem im unberührten Rückzugsgebiet
unterscheiden kann.

\paragraph{Qualität der Zeitstempel.} Die Zeitstempel stammen aus der automatisierten
Auswertung der in die Aufnahmen eingeblendeten Textzeile. Fehlerhaft erkannte Gerätekennungen
sind in der Auswertung ausgeschlossen worden; eine entsprechende Prüfung der Zeitangaben
selbst fand nicht statt. Ein Teil der Zeitstempel kann daher fehlerhaft sein, ohne dass dies
auffiele.

\paragraph{Mehrfach vorliegende Ortungen.} Die in Abschnitt~\ref{subsec:datenhaltung}
beschriebene Doppelung der Karpaten-Ortungen betrifft die Darstellung und abgeleitete
Kennzahlen. Für den hier ausgewerteten Vergleich ist sie ohne Belang, da dieser auf den
Kameradaten und der Geschwindigkeitstabelle beruht, nicht auf der Ortungstabelle. Für
künftige Auswertungen auf dieser Tabelle ist sie jedoch zu beheben.

\subsection{Modellannahmen}
\label{subsec:limit-modell}

\paragraph{Artunterschied.} Die japanischen Kameras stehen im Verbreitungsgebiet des
Asiatischen Schwarzbären, während das Bewegungsmodell aus Daten des Europäischen Braunbären
stammt. Das ist die weitreichendste inhaltliche Vereinfachung dieser Arbeit. Sie ist bewusst
getroffen, weil für die japanische Art keine vergleichbaren Trajektoriendaten vorlagen, und
sie wird ausdrücklich nicht als Gleichsetzung der beiden Arten verstanden. Jede Aussage über
die japanischen Verhältnisse steht unter diesem Vorbehalt.

\paragraph{Übertragene Parameter.} Dasselbe gilt für die Geschwindigkeitstabelle und die
gemessene Richtungstreue: Beide stammen aus den Karpaten. Die Richtungstreue ist zudem, wie
in Abschnitt~\ref{subsec:richtungstreue} dargelegt, bei kurzen Schritten vom Messfehler der
GPS-Ortung überlagert, sodass die dort angegebenen negativen Werte als Obergrenze der
tatsächlichen Gewundenheit zu lesen sind.

\paragraph{Vereinfachtes Geländemodell.} Die Gewichte der Geländebewertung sind von Hand
gesetzte Näherungen und keine aus Daten geschätzten Koeffizienten. Eine integrierte
Step-Selection-Analyse \cite{Avgar2016} würde diese Werte aus beobachteten Schritten
schätzen, statt sie vorzugeben, und ist damit der methodisch saubere Weg. Er setzt jedoch
Trajektoriendaten aus dem betrachteten Gebiet voraus, die für Japan gerade fehlen.

Hinzu kommt eine räumliche Grenze: Das lokal vorgehaltene Gelände umfasst die Kamerastandorte
zuzüglich eines Puffers. Bei Simulationsläufen über 48~Stunden verlassen Bären dieses Gebiet,
und die Geländegewichtung entfällt dort stillschweigend. Für belastbare Läufe über 24~Stunden
hinaus müsste der Puffer vergrößert werden.

\subsection{Technische Grenzen des Systems}
\label{subsec:technische-grenzen}

\paragraph{Zeitstempel als Zeichenketten.} Die Zeitstempel liegen weiterhin als Zeichenketten
in der Datenbank. Eine Umstellung auf einen echten Zeitstempeltyp hätte eine Datenmigration
und die Anpassung aller darauf aufsetzenden Abfragen erfordert; der Nutzen wäre vor allem bei
deutlich größeren Datenmengen spürbar geworden. Die Einschränkung bleibt bestehen: Eine
Filterung nach Zeiträumen ist auf der Serverseite nur eingeschränkt wirksam.

\paragraph{Mehrfach vorliegende Ortungen.} Die Abfragen auf der Ortungstabelle filtern die
Doppelungen nicht heraus. Praktisch wirkt sich das derzeit kaum aus, weil dieselbe Position
zweimal übereinander gezeichnet wird und optisch nicht auffällt. Bedeutsam würde es, sobald
aus dieser Tabelle Kennzahlen abgeleitet werden - etwa Schrittlängen oder Tagesdistanzen --,
da dann Schritte der Länge null in die Statistik eingingen.

\paragraph{Abhängigkeit von externen Diensten.} Kartenkacheln, Höhenmodell und die beiden
Vergleichsdatenquellen stammen von Diensten außerhalb der Hochschule. Jeder davon kann
ausfallen oder seine Nutzungsbedingungen ändern. Das System ist darauf vorbereitet, da ein
Ausfall jeweils zu einem Hinweis und nicht zu einem Abbruch führt, aber die betroffene
Funktion steht dann nicht zur Verfügung.

\paragraph{Keine Mehrbenutzerabsicherung.} Mehrere Personen können gleichzeitig mit dem
System arbeiten, da die Daten zentral liegen. Schreibende Zugriffe - das Anlegen und Löschen
von Markern und Flächen - sind jedoch nicht gegen gleichzeitige Bearbeitung abgesichert.
Bei der zu erwartenden Zahl gleichzeitiger Nutzender in einer Forschungsgruppe ist das
vertretbar, wäre bei breiterem Einsatz aber nachzuholen.

% ====================================================================
\chapter{Fazit und Ausblick}
\label{chap:fazit}

\section{Zusammenfassung der Ergebnisse}
\label{sec:zusammenfassung}

Ziel dieser Arbeit war, drei getrennt entstandene Teilsysteme der rechnergestützten
Naturbeobachtung in einer gemeinsamen Kartenplattform zusammenzuführen und darüber hinaus
Werkzeuge bereitzustellen, mit denen sich simulierte Bewegungen gegen reale Beobachtungsdaten
prüfen lassen. Beides wurde erreicht.

\paragraph{Integration und technische Basis.} An die Stelle zweier voneinander unabhängiger
Karteninstanzen ist eine gemeinsame, auf MapLibre~GL~JS beruhende Kartenkomponente getreten,
auf der alle Datenbestände gleichzeitig liegen und einzeln ein- und ausgeblendet werden
können. Das Bewegungsmodell, das zuvor ausschließlich als eigenständiger Quelltext ohne
Oberfläche vorlag, ist in die Anwendung eingebunden; die zuvor auf zwei Tabellen verteilten
Kameraaufnahmen sind in einer zusammengeführt, und die Zugangsdaten sind aus dem Quelltext
ausgelagert. Ergänzt wurden zwei unabhängige öffentliche Datenquellen als Vergleichsebene.

Zwei Entwurfsentscheidungen prägen das Ergebnis. Die Karte ist als statische Komponente ohne
Rückkanal eingebunden, was die Interaktion schnell macht und den Betrieb ohne zusätzliche
Übersetzungsschritte ermöglicht; der fehlende Rückkanal wird für die wenigen schreibenden
Vorgänge über die Adresszeile umgangen. Und die Schrittauswahl des Bewegungsmodells wurde
neu geschrieben statt erweitert, weil die Vorarbeit zwar Kandidaten erzeugte, unter ihnen
aber keine Auswahl traf.

\paragraph{Analyse- und Validierungswerkzeuge.} Der zeitliche Ablauf einer Bewegung ist
abspielbar geworden, und zwar in synchroner Gegenüberstellung von realer Trajektorie und
Modellausgabe für dasselbe Tier, denselben Tag und dieselben Zeitpunkte. Die Richtungswahl
des Modells wurde auf drei Wegen verbessert: durch eine aus 46.161~realen Schrittpaaren
geschätzte, nach Schrittlänge aufgeschlüsselte Richtungstreue, durch eine Geländegewichtung
anhand lokal vorgehaltener OpenStreetMap-Daten und durch die Beseitigung einer systematischen
Winkelverzerrung, die zuvor ein Drittel aller Richtungen unerreichbar machte.

Für die Übertragung auf die japanischen Kameradaten wurde ein Prüfverfahren entwickelt, das
dem Datentyp angemessen ist: Da sich aus punktuellen Detektionen keine Pfade rekonstruieren
lassen, wird nicht räumlich, sondern zeitlich verglichen - über den Überlappungskoeffizienten
der Tagesaktivitätsverteilungen, mit einer Unsicherheitsangabe aus einem Bootstrap-Verfahren.

\paragraph{Befunde zu den Vorarbeiten.} Im Verlauf der Analyse traten mehrere Eigenschaften
der übernommenen Arbeiten zutage, die für den weiteren Betrieb von Bedeutung sind: die
fehlende Auswahl zwischen den erzeugten Kandidatenschritten, die auf eine Stunde festgelegte
Schrittlänge bei tatsächlich ungleichmäßigen Ortungsabständen, die doppelt geführten
Bildtabellen und die mehrfach vorliegenden Ortungen. Sie sind in
Kapitel~\ref{chap:analyse} und~\ref{chap:diskussion} jeweils dort benannt, wo sie
Entwurfsentscheidungen begründen.

\paragraph{Einordnung.} Das System ist ein Arbeitswerkzeug für ein Forschungsteam, kein
Vorhersagewerkzeug. Die simulierten Pfade sind Zufallsziehungen aus einem Modell, und das
ermittelte Übereinstimmungsmaß ist angesichts des Stichprobenumfangs ein explorativer Befund.
Der Beitrag dieser Arbeit liegt darin, dass sich diese Frage überhaupt stellen und
quantitativ beantworten lässt - im Ausgangszustand war weder das eine noch das andere
möglich.

\section{Ausblick auf zukünftige Erweiterungen}
\label{sec:ausblick}

\paragraph{Bewegungsparameter der japanischen Art.} Die weitreichendste Vereinfachung dieser
Arbeit ist die Übertragung von Parametern des Europäischen Braunbären auf ein Gebiet des
Asiatischen Schwarzbären. Liegen für die japanische Art veröffentlichte Werte zu
Geschwindigkeit und Richtungsverhalten vor, ließen sie sich unmittelbar einsetzen - die
Geschwindigkeitstabelle ist nach Art gegliedert und müsste lediglich um einen Eintrag
ergänzt werden. Der Vergleich beider Parametersätze gegen dieselben Kameradaten wäre für
sich genommen ein lohnender Untersuchungsgegenstand.

\paragraph{Datengestützte Geländegewichtung.} Die Gewichte der Geländebewertung sind von Hand
gesetzt. Stünden Trajektoriendaten aus dem japanischen Gebiet zur Verfügung, ließen sich die
Koeffizienten über eine integrierte Step-Selection-Analyse \cite{Avgar2016} aus beobachteten
Schritten schätzen; dafür bestehen eingeführte Werkzeuge \cite{Signer2019}. Damit entfiele
die derzeit schwächste Stelle der Modellerweiterung.

\paragraph{Erweiterung des vorgehaltenen Geländes.} Unabhängig davon ist das lokal
vorgehaltene Gebiet zu vergrößern. Derzeit umfasst es die Kamerastandorte zuzüglich eines
Puffers, den Simulationsläufe über 48~Stunden verlassen. Der Vorgang ist bereits als Werkzeug
umgesetzt und muss lediglich mit einem größeren Puffer erneut ausgeführt werden.

\paragraph{Kameraindividuelle Auswertung.} Wächst der Bestand an Aufnahmen, wird eine
Auswertung je Kamera möglich. Interessant wäre insbesondere der Vergleich zwischen Kameras in
Waldlage und in Siedlungsnähe: Unterscheiden sich deren Tagesrhythmen, wäre das ein Hinweis
auf eine Verhaltensanpassung an menschliche Anwesenheit - eine Frage, die mit der
vorliegenden Datenmenge nicht beantwortet werden kann.

\paragraph{Technische Anschlusspunkte.} Drei Punkte sind unabhängig von den inhaltlichen
Fragen zu erledigen: die Umstellung der Zeitstempel auf einen echten Zeitstempeltyp mitsamt
serverseitiger Zeitraumfilterung, die Herausfilterung der mehrfach vorliegenden Ortungen und
das Entfernen der Rückfallwerte für die Zugangsdaten aus dem Quelltext. Der dritte Punkt ist
vor einem produktiven Einsatz zwingend, zusammen mit dem Wechsel des in der Versionsgeschichte
enthaltenen Zugangsworts.

% ====================================================================
\printbibliography[title={Literatur}]

\end{document}
